% Εσωτερική έκδοση v16.5 (28/10/2026). Ελεγχόμενη αναθεώρηση της v16.4, με B2--EISE και συντακτική εναρμόνιση.
% Μεταγλώττιση: XeLaTeX, δύο περάσματα. Δεν απαιτούνται εξωτερικά αρχεία γραμματοσειρών.
\documentclass[a4paper]{article}
\usepackage[a4paper,left=1.9cm,right=1.9cm,top=1.7cm,bottom=1.75cm,headheight=14pt,headsep=6mm]{geometry}
\usepackage{fix-cm}
\usepackage{fontspec,polyglossia}
\setdefaultlanguage{greek}
\setotherlanguage{english}
% Φορητή επιλογή: GNU FreeFont, αλλιώς Liberation (χρησιμοποιήθηκε στον παρόντα έλεγχο).
\IfFontExistsTF{FreeSerif}{\def\MainFace{FreeSerif}\def\SansFace{FreeSans}\def\MonoFace{FreeMono}}{\def\MainFace{Liberation Serif}\def\SansFace{Liberation Sans}\def\MonoFace{Liberation Mono}}
\setmainfont{\MainFace}[Ligatures=TeX]
\newfontfamily\greekfont{\MainFace}[Ligatures=TeX]
\newfontfamily\englishfont{\MainFace}[Ligatures=TeX]
\setsansfont{\SansFace}[Scale=MatchLowercase]
\newfontfamily\greekfontsf{\SansFace}[Scale=MatchLowercase]
\setmonofont{\MonoFace}[Scale=MatchLowercase]
\newfontfamily\greekfonttt{\MonoFace}[Scale=MatchLowercase]
\usepackage{microtype,amsmath,amssymb,graphicx,tabularx,booktabs,array,ragged2e,calc}
\usepackage{xcolor,fancyhdr,titlesec,enumitem,needspace}
\usepackage[font=footnotesize,labelfont=bf,labelsep=period,justification=justified,singlelinecheck=false,skip=3pt]{caption}
\usepackage{tikz}
\usetikzlibrary{positioning,arrows.meta,shapes.geometric,calc}
\usepackage{indentfirst}
\usepackage{xurl}
\usepackage{hyperref}
\urlstyle{same}
\definecolor{acc}{HTML}{7A2E2E}
\definecolor{soft}{HTML}{F7F1F1}
\definecolor{edge}{HTML}{8A8A8A}
\definecolor{routeb}{HTML}{2E4F7A}
\definecolor{softb}{HTML}{EEF2F8}
\definecolor{ecsi}{HTML}{4F6B2E}
\definecolor{softe}{HTML}{F2F6EC}
\hypersetup{colorlinks=true,linkcolor=acc,citecolor=acc,urlcolor=acc,
 pdftitle={Η μουσική θέση ως υπολογιστική μονάδα της βυζαντινής μουσικής — σύντομη ερευνητική πρόταση},
 pdfauthor={Γεώργιος Σπυριδάκης},
 pdfsubject={Μουσικολογικά ενημερωμένη μηχανική μάθηση, συμβολική ανάκτηση μουσικής πληροφορίας και βυζαντινή μουσικολογία}}
\newcolumntype{Y}{>{\raggedright\arraybackslash}X}
\newcolumntype{L}[1]{>{\raggedright\arraybackslash}p{#1}}
\setlength{\parindent}{1em}
\setlength{\parskip}{0.6pt}
\setlength{\tabcolsep}{3.5pt}
\renewcommand{\arraystretch}{1.1}
\emergencystretch=2.5em
\raggedbottom
\widowpenalty=10000
\clubpenalty=10000
\titleformat{\section}{\large\bfseries\color{acc}}{\thesection.}{0.5em}{}
\titleformat{\subsection}{\normalsize\bfseries\color{acc}}{\thesubsection}{0.5em}{}
\titleformat{\paragraph}[runin]{\normalsize\bfseries}{}{0pt}{}[.]
\titlespacing*{\section}{0pt}{6pt}{2.5pt}
\titlespacing*{\subsection}{0pt}{5pt}{2pt}
\titlespacing*{\paragraph}{0pt}{2pt}{0.5em}
\newcommand{\sectionbreak}{\Needspace{6\baselineskip}}
\newcommand{\subsectionbreak}{\Needspace{5\baselineskip}}
\pagestyle{fancy}
\fancyhf{}
\lhead{\footnotesize Γ. Σπυριδάκης / Σύντομη ερευνητική πρόταση}
\rhead{\footnotesize\thepage}
\renewcommand{\headrulewidth}{0.3pt}
\captionsetup[table]{position=top}
\newcommand{\keybox}[1]{\par\smallskip\noindent\fcolorbox{ecsi}{softe}{\parbox{\dimexpr\linewidth-2\fboxsep-2\fboxrule\relax}{\setlength{\parindent}{0pt}\setlength{\parskip}{2pt}#1}}\par\smallskip}
% Σύμβολα
\newcommand{\CE}{\mathrm{CE}}
\newcommand{\Mth}{\ensuremath{M_{\text{θέσις}}}}
\newcommand{\dhist}{\ensuremath{\Delta_{\mathrm{hist}}}}
\newcommand{\dlearn}{\ensuremath{\Delta_{\mathrm{learn}}}}
\newcommand{\drec}{\ensuremath{\Delta_{\mathrm{rec}}}}
\newcommand{\dtwo}{\ensuremath{\Delta_{\mathrm{2ex}}}}
\newcommand{\drule}{\ensuremath{\Delta^{\mathrm{rule}}_{\mathrm{rec}}}}
\newcommand{\dfam}{\ensuremath{\Delta_{\mathrm{family}}}}
\newcommand{\datl}{\ensuremath{\Delta_{\mathrm{ATL}}}}
\newcommand{\Rs}{\ensuremath{R^{\star}}}
\newcommand{\MAP}{\ensuremath{\mathrm{MacroAP}_{\mathrm{hard}}}}
\newcommand{\Qreal}{\ensuremath{Q_{\mathrm{real}}}}
\newcommand{\qpar}[1]{\par\smallskip\noindent\textbf{#1}\enspace}
\newenvironment{reflist}{\begin{list}{}{\leftmargin=1.3em\itemindent=-1.3em\itemsep=0pt\parsep=0pt}\fontsize{8}{8.7}\selectfont}{\end{list}}
\tikzset{
 box/.style={draw=edge,rounded corners=2.5pt,line width=0.6pt,align=center,font=\footnotesize,inner sep=3pt},
 boxa/.style={box,draw=acc,fill=soft},
 boxb/.style={box,draw=routeb,fill=softb},
 boxe/.style={box,draw=ecsi,fill=softe},
 arr/.style={-{Stealth[length=2mm]},draw=edge,line width=0.6pt},
}
\begin{document}
\fontsize{10.5}{12.6}\selectfont
\setlength{\parindent}{1em}

% ------------------------------------------------------------------ ΕΞΩΦΥΛΛΟ
\hypersetup{pageanchor=false}
\begin{titlepage}
\centering
\vspace*{0.8cm}
{\sffamily\large ΕΛΛΗΝΙΚΟ ΜΕΣΟΓΕΙΑΚΟ ΠΑΝΕΠΙΣΤΗΜΙΟ\par}
{\sffamily Τμήμα Μουσικής Τεχνολογίας και Ακουστικής\par}
\vspace{1.7cm}
{\Large\bfseries\color{acc} ΕΡΕΥΝΗΤΙΚΗ ΠΡΟΤΑΣΗ ΔΙΔΑΚΤΟΡΙΚΗΣ ΔΙΑΤΡΙΒΗΣ\par}
\vspace{0.25cm}
{\large Σύντομη έκδοση\par}
\vspace{1.1cm}
{\fontsize{22}{27}\selectfont\bfseries Η μουσική θέση ως υπολογιστική μονάδα\\[2pt] της βυζαντινής μουσικής\par}
\vspace{0.6cm}
{\Large Συμβολικό MIR, μάθηση αναπαραστάσεων και μουσικολογικά ενημερωμένη μηχανική μάθηση\par}
\vspace{0.5cm}
{\large\itshape The musical thesis as a computational unit of Byzantine chant:\\ symbolic MIR, representation learning and musicology-informed machine learning\par}
\vfill
\begin{tabularx}{0.94\linewidth}{@{}>{\bfseries}L{3.9cm}Y@{}}
Υποψήφιος & Γεώργιος Σπυριδάκης\\[7pt]
Γνωστικό πεδίο & Συμβολικό Music Information Retrieval (MIR), μάθηση αναπαραστάσεων και μουσικολογικά ενημερωμένη μηχανική μάθηση\\[7pt]
Μουσικό πεδίο & Μουσικές θέσεις στα παπαδικά Χερουβικά και Κοινωνικά της Νέας Μεθόδου\\[7pt]
Χρονικός ορίζοντας & 48 μήνες\\[7pt]
Τεκμηρίωση & Η πρόταση συνοδεύεται από τεχνικά πρωτόκολλα Α2 και Μ2/B2--EISE. Αναλυτικός σχεδιασμός και μουσικολογική τεκμηρίωση στο \href{https://github.com/GeorgeSpy/byz-position-phd-study}{github.com/GeorgeSpy/byz-position-phd-study}\\[7pt]
Ημερομηνία & Οκτώβριος 2026
\end{tabularx}
\vspace{0.4cm}
\end{titlepage}
\setcounter{page}{1}
\hypersetup{pageanchor=true}

% ============================================================= ΠΕΡΙΛΗΨΗ
\section*{Περίληψη}
Η βυζαντινή ψαλτική είναι μονόφωνη φωνητική μουσική της ορθόδοξης λατρείας, που μεταδόθηκε επί αιώνες με τη γραφή και την προφορική διδασκαλία μαζί. Η σημειογραφία της δηλώνει την κίνηση της φωνής από φθόγγο σε φθόγγο, και στις παλαιότερες μορφές της συνδυάζει τα σημάδια των διαστημάτων με μεγάλα σημάδια χαρακτηριστικής μελικής ενέργειας. Η σχέση αυτής της γραφής με την αναλυτική απόδοση της Νέας Μεθόδου του 1814 είναι κεντρικό ζήτημα της βυζαντινής μουσικολογίας, και οι εξηγήσεις του Γρηγορίου Πρωτοψάλτη και του Χουρμουζίου Χαρτοφύλακος, που μετέφεραν παλαιότερα μέλη στη νέα γραφή, σώζουν για ορισμένα χωρία δύο αποδόσεις του ίδιου μουσικού κειμένου. Σε όλη αυτή τη διαδρομή το μέλος οργανώνεται με ένα λεξιλόγιο παραδοσιακών σχηματισμών, τις μουσικές θέσεις, που επανέρχονται από έργο σε έργο προσαρμοσμένες στον ήχο, δηλαδή στον τρόπο του μέλους, στη βαθμίδα και στο κείμενο, και που ήδη η θεωρία του 15ου αιώνα παρομοίαζε με λέξεις.

Ως λέξεις της μουσικής γλώσσας, οι θέσεις βρίσκονται ανάμεσα στο μεμονωμένο σημάδι και στο πλήρες μέλος, εκεί όπου η γραφή αποκτά μουσικό νόημα, και γι' αυτό τις επιλέγουμε ως μονάδα της έρευνας. Η ταυτότητά τους όμως δεν συμπίπτει με την ομοιότητα: συγγενικές μελωδικές πορείες μπορεί να ανήκουν σε διαφορετικές θέσεις, η ίδια θέση αναπτύσσεται με διαφορετική έκταση και μορφή, και τα όριά της δεν είναι πάντοτε μονοσήμαντα. Η απόσταση αυτή κάνει το μουσικολογικό ερώτημα απαιτητικό πρόβλημα συμβολικής ανάκτησης μουσικής πληροφορίας (Music Information Retrieval, MIR) και μηχανικής μάθησης, αφού ένα σύστημα πρέπει να εντοπίζει στη συνεχή γραφή μονάδες μεταβλητής έκτασης, να αποφασίζει σε ποια οικογένεια θέσεων ανήκουν και να απέχει όταν τα τεκμήρια δεν αρκούν. Η διατριβή το αντιμετωπίζει ως μουσικολογικά ενημερωμένη μηχανική μάθηση (Musicology-Informed Machine Learning, MIML): η μουσικολογία ορίζει τη μονάδα και τις οικογένειες με προδιαγραφές που ελέγχουν ανεξάρτητοι ειδικοί, δίνει τις σχέσεις και τους κανόνες από τους οποίους μαθαίνει το μοντέλο, και γίνεται η ίδια πειραματική μεταβλητή, ώστε να μετριέται τι προσφέρει. Ο κύριος μεθοδολογικός στόχος ανήκει στο συμβολικό MIR και στη μηχανική μάθηση, ενώ ο ελέγξιμος ορισμός της θέσης και των οικογενειών της συνιστά συγχρόνως ουσιαστικό μουσικολογικό ερευνητικό αποτέλεσμα. Το πλαίσιο αναπτύσσεται στα παπαδικά Χερουβικά και Κοινωνικά της Νέας Μεθόδου, αργά και μελισματικά μέλη της Θείας Λειτουργίας, όπου οι θέσεις είναι εκτενείς και οι δύσκολες περιπτώσεις πολλές.

Από αυτή την αφετηρία η έρευνα ακολουθεί τρία νήματα, που στηρίζονται στα ίδια πραγματικά, επισημειωμένα έργα. Το πρώτο αφορά την ίδια την αναγνώριση: ποια ανάγνωση της γραφής και ποια στρώματα πληροφορίας, σημειογραφικά, χρονικά, τροπικά και κειμενικά, επιτρέπουν στο μοντέλο να οριοθετεί και να αναγνωρίζει τις θέσεις μαζί, και πόσα από τα όρια εξηγούν ήδη οι γενικές ενδείξεις μελωδικής κατάτμησης (Μελέτη 1, Μ1). Το δεύτερο εξετάζει την ίδια τη μουσικολογική γνώση. Οι προδιαγραφές των οικογενειών έχουν λογική δομή, με απαιτούμενα, εναλλακτικά και αποκλειόμενα γνωρίσματα· το πείραμα συγκρίνει έγκυρες με ελεγχόμενα διαταραγμένες προδιαγραφές, σε επίπεδη ή δομημένη χρήση, και εκτιμά την επίδραση του έγκυρου περιεχομένου και το πρόσθετο όφελος της λογικής σύνθεσής του (πείραμα Α2). Το τρίτο ρωτά τι μαθαίνει ένας τμηματικός κωδικοποιητής, δηλαδή ένα νευρωνικό μοντέλο που μετατρέπει κάθε τμήμα της γραφής σε αριθμητική αναπαράσταση, όταν παρατηρεί πολλές πραγματώσεις της ίδιας οικογένειας, και αν κερδίζει σε αναγνώριση χωρίς να σβήνει τη μουσικά σημαντική παραλλαγή (τρεις συγκρίσιμες συνθήκες εκπαίδευσης, B0--B2). Η Μ2 περιλαμβάνει επίσης την απογραφή και τεκμηρίωση συν-εξηγήσεων, δύο ιστορικών αποδόσεων του ίδιου χωρίου. Όπου υπάρχουν επιλέξιμα ζεύγη, εξετάζει πιλοτικά αν αυτή η εποπτεία οργανώνει διαφορετικά την παραλλαγή από την κοινή οικογένεια· η επιβεβαιωτική σύγκριση απαιτεί επιπλέον ανεξάρτητες εκτελέσεις και ελεγμένη στατιστική επάρκεια.

Ένας πιλότος των πρώτων έξι μηνών διατρέχει όλη την αλυσίδα, από την πηγή ως το πρώτο μοντέλο, και τα προκαθορισμένα κριτήρια επάρκειας καθορίζουν την κλίμακα κάθε ανάλυσης πριν φανούν τα αποτελέσματα. Στο τέλος αυτής της διαδρομής η διατριβή επιδιώκει να παραδώσει μια νέα εργασία συμβολικού MIR με σώμα αξιολόγησης, ελεγμένες προδιαγραφές οικογενειών θέσεων και ελεγχόμενη εκτίμηση του τι προσθέτει η μουσικολογική γνώση στη μηχανική μάθηση. Επειδή μελωδικοί τύποι και συνοπτικές γραφές υπάρχουν και σε άλλες παραδόσεις, το πλαίσιο σχεδιάζεται ώστε να δοκιμαστεί και εκεί, με πρώτο ορίζοντα το ζναμένιο μέλος, την ιστορικά συγγενή νευματική παράδοση της ρωσικής εκκλησίας.

\smallskip\noindent\textbf{Λέξεις-κλειδιά:} βυζαντινή ψαλτική, μουσική θέση, συμβολικό MIR, μουσικολογικά ενημερωμένη μηχανική μάθηση, κοινή οριοθέτηση και αναγνώριση, μάθηση αναπαραστάσεων, εξήγηση.

% ============================================================= 1
\section{Το μουσικό πρόβλημα: ταυτότητα μέσα στην παραλλαγή}\label{sec:intro}
\subsection{Η θέση ως υπολογιστική μονάδα}\label{sec:unit}
Η απόσταση ανάμεσα σε ταυτότητα και ομοιότητα γίνεται συγκεκριμένη μόλις διαβάσει κανείς ένα παπαδικό μέλος της Νέας Μεθόδου. Η γραφή αποδίδει τη μελωδία με ποσοτικά σημεία, που δηλώνουν σχετική διαστηματική κίνηση, συνοδευμένα από χρονικές και ποιοτικές ενδείξεις, και διαβάζεται πάντοτε μέσα στο τροπικό πλαίσιο του ήχου. Η \emph{μετροφωνία}, η ιστορική πρακτική να διαβάζεται η διαστηματική ποσότητα των σημείων ως διαδοχή βημάτων, δίνει τον διαστηματικό σκελετό του μέλους, αλλά δεν εξαντλεί την ταυτότητα μιας θέσης. Στη μορφολογική του μελέτη των Χερουβικών, ο Κ.~Χ.~Καραγκούνης αναζητά κοινές θέσεις σε έργα διαφορετικών μελουργών χωρίς να στηρίζεται αποκλειστικά στη μετροφωνία ή στη βαθμίδα (2003, σ.~783), και διακρίνει τη θέση του αντικενώματος από σχηματισμούς με την ίδια μετροφωνία που η παλαιότερη γραφή συνδέει με το λύγισμα ή με το έτερον παρακάλεσμα (σ.~788). Τη διαφορά τη φέρουν εκεί οι \emph{μεγάλες υποστάσεις}, τα ιστορικά σημάδια που δηλώνουν χαρακτηριστική μελική ενέργεια πέρα από τη διαστηματική πορεία, και έτσι σχηματισμοί με συγγενή επιφάνεια αποκτούν διαφορετικές ταυτότητες.

Το αντίστροφο φαινόμενο είναι εξίσου διδακτικό, γιατί η ίδια θέση μπορεί να αλλάζει μορφή για να προσαρμοστεί στο περιβάλλον της. Στο Χερουβικό του μελουργού Μανουήλ Αγαλλιανού η εναρκτήρια θέση διαφοροποιείται ανάλογα με τη σημειωμένη \emph{μαρτυρία}, την ένδειξη που δηλώνει την τροπική βαθμίδα, ώστε να μένει σταθερό το τονικό σημείο εκκίνησης του μέλους (Καραγκούνης, 2005, σ.~175)· η μαρτυρία αυτή διακρίνεται από τον ιστορικό \emph{μάρτυρα}, δηλαδή το χειρόγραφο ή το έντυπο που διασώζει τη γραφή. Με τον ίδιο τρόπο, ο Ι.~Αρβανίτης δείχνει ότι μια θέση μένει αναγνωρίσιμη όταν αλλάζουν ο αριθμός των συλλαβών, ο λεκτικός τονισμός ή η σημειογραφική της απόδοση (2010, τ.~Α΄, σσ.~136--142). Η επιφανειακή ομοιότητα μπορεί λοιπόν να είναι μεγάλη ενώ η ταυτότητα διαφέρει, και μικρή ενώ η ταυτότητα μένει η ίδια.

Για να μελετηθεί αυτή η ασυμμετρία με ακρίβεια, η διατριβή ξεχωρίζει έννοιες που στην καθημερινή χρήση συγχέονται. Όταν ένας ειδικός αναγνωρίζει σε ένα χωρίο Χερουβικού τη θέση που η παράδοση ονομάζει παρακλητική, το συγκεκριμένο χωρίο στη συγκεκριμένη πηγή είναι μια \emph{εμφάνιση} της θέσης, και η μελική μορφή που αποκτά εκεί, με τον ήχο, τη βαθμίδα και τη συλλαβική της κατανομή, είναι η \emph{πραγμάτωσή} της. Εμφανίσεις από διαφορετικά έργα που αναγνωρίζονται ως εκδοχές του ίδιου μελικού τύπου ανήκουν στην ίδια \emph{οικογένεια θέσεων}, και την οικογένεια την ορίζει μια γραπτή \emph{προδιαγραφή}: καταγράφει τις απαιτούμενες και τις υποστηρικτικές μουσικές σχέσεις, τους εναλλακτικούς συνδυασμούς που αρκούν για την ένταξη, τους τροπικούς όρους και τα κριτήρια αποκλεισμού, σε αριθμημένες εκδόσεις που άλλοι ειδικοί μπορούν να εφαρμόσουν σε νέο υλικό. Η ιστορική ονομασία, εδώ «παρακλητική», είναι πολύτιμο τεκμήριο αλλά όχι αυτόματος ορισμός, αφού μια ονομασία μπορεί να αντιστοιχεί σε μία οικογένεια, να μοιράζεται σε περισσότερες ή να συγχωνεύεται με άλλη. Η εμφάνιση είναι λοιπόν αυτό που εντοπίζεται μέσα στο έργο, ενώ η οικογένεια είναι αυτό που αναγνωρίζεται σε αυτήν. Η μετατροπή αυτής της διάκρισης σε υπολογιστικά ελέγξιμο ορισμό δεν αποτελεί απλή ονοματοδοσία ή επισήμανση δεδομένων: προϋποθέτει αιτιολογημένη επιλογή γνωρισμάτων, προσδιορισμό των αποδεκτών παραλλαγών και τεκμηρίωση των περιπτώσεων όπου η ένταξη παραμένει αμφίσημη.

Η εμφάνιση έχει μεταβλητό και κάποτε αμφίσημο μήκος, που δεν ακολουθεί μηχανικά τη συλλαβή. Μια εναρκτήρια θέση μπορεί να καλύπτει τις τρεις πρώτες συλλαβές «Οἱ τὰ Χε-», ενώ αλλού μία μόνο συλλαβή αναπτύσσεται σε περισσότερες διαδοχικές θέσεις (Καραγκούνης, 2003, σσ.~207, 209), και βραχείς σχηματισμοί με συνδετική λειτουργία παρεμβάλλονται ανάμεσα σε εκτενέστερες ακολουθίες (σ.~192). Επειδή η αναγνώριση εξαρτάται από την έκταση του τμήματος και η έκταση από το τι αναγνωρίζεται, η διατριβή εκτιμά όρια και οικογένεια μαζί. Η μονάδα τοποθετείται έτσι ανάμεσα στο μεμονωμένο σημείο και στην ακολουθία των θέσεων που οργανώνει το πλήρες έργο, και αυτή η ενδιάμεση κλίμακα είναι το αντικείμενο του κορμού της έρευνας.

\subsection{Η εξήγηση: μια σπάνια ιστορική σχέση}\label{sec:exegesis}
Η ψαλτική διαθέτει μια ιδιαιτερότητα που σπάνια συναντά κανείς σε άλλες παραδόσεις. Η \emph{εξήγηση}, η αναλυτικότερη απόδοση παλαιότερων μελών στη Νέα Μέθοδο (Γ.~Θ.~Στάθης, 1978), διασώζει για ορισμένα χωρία δύο αποδόσεις του ίδιου παλαιότερου μέλους, του Γρηγορίου και του Χουρμουζίου. Η σχέση των δύο αποδόσεων δεν προκύπτει από κάποιο μέτρο ομοιότητας· διαπιστώνεται από την κριτική εξέταση των πηγών, δηλαδή από το ποιο χωρίο εξηγείται, πώς αντιστοιχούν οι συλλαβές και ποια εξάρτηση αντιγραφής ή κοινής πηγής ενδέχεται να συνδέει τους δύο εξηγητές. Δύο τέτοιες \emph{συν-εξηγήσεις} δείχνουν ποια παραλλαγή είναι συμβατή με την ίδια θέση στο ίδιο ιστορικό χωρίο, και αποτελούν γι' αυτό μια στενότερη σχέση από την απλή κοινή οικογένεια. Το πόσα τέτοια ζεύγη επαληθεύονται είναι όμως άγνωστο πριν από την απογραφή. Γι’ αυτό η Μ2 οργανώνεται σε δύο σκέλη: τη μάθηση από εμφανίσεις κοινής οικογένειας (Μ2a) και την ιστορικά τεκμηριωμένη εποπτεία από συν-εξηγήσεις (Μ2b). Η απογραφή της Μ2b αρχίζει στο πρώτο εξάμηνο, ενώ η πιλοτική εκπαίδευση και η επιβεβαιωτική αξιολόγηση εξαρτώνται από την πραγματική επάρκεια πηγών και ανεξάρτητων εκτελέσεων (§\ref{sec:historicpriority}).

\subsection{Η μουσικολογική γνώση ως πειραματική μεταβλητή}\label{sec:miml}
Τόσο οι προδιαγραφές των οικογενειών όσο και οι συν-εξηγήσεις είναι μουσικολογική γνώση που μπορεί να καθοδηγήσει τη μάθηση ενός μοντέλου. Η διατριβή όμως δεν τη χρησιμοποιεί μόνο ως πηγή ετικετών· την αντιμετωπίζει ως πειραματική μεταβλητή. Η πληροφορημένη μηχανική μάθηση (informed machine learning) περιγράφει πώς μια προγενέστερη γνώση εισέρχεται ρητά στα δεδομένα, στη μορφή του μοντέλου, στη συνάρτηση που βελτιστοποιείται ή στον έλεγχο του αποτελέσματος (von~Rueden κ.ά., 2023· Bauckhage, Schulz και von~Rueden, 2025), και στο MIR τα μοντέλα που ενσωματώνουν μουσική γνώση αποτελούν πια χωριστή κατεύθυνση (Richard κ.ά., 2024). Στην ιστορική μουσικολογία, όμως, η γνώση δεν είναι βέβαιη όπως ένας φυσικός νόμος, γιατί οι οικογένειες θέσεων είναι αναθεωρήσιμες υποθέσεις. Γι' αυτό ο όρος MIML χρησιμοποιείται εδώ περιγραφικά, για ένα πλαίσιο με δύο γνωρίσματα: κάθε ρόλος της γνώσης συνοδεύεται από δικό του έλεγχο (Πίνακας~\ref{tab:knowledge}), και η ίδια η γνώση μεταβάλλεται πειραματικά, ώστε να μετριέται αν ωφελεί το περιεχόμενό της ή μόνο η μορφή της. Αυτή η σχέση προσδιορίζει και το αντικείμενο της διατριβής: η κύρια ερευνητική εργασία είναι ο σχεδιασμός και η συγκριτική αξιολόγηση υπολογιστικών μεθόδων συμβολικού MIR και μηχανικής μάθησης. Η ιστορική και μουσικολογική τεκμηρίωση καθορίζει το μουσικό αντικείμενο, θεμελιώνει τις διαγνωστικές σχέσεις, παρέχει επισημάνσεις και ---όπου υπάρχουν επαληθευμένες αντιστοιχίες--- σχεσιακή εποπτεία, και προσφέρει ανεξάρτητα κριτήρια μουσικής εγκυρότητας. Δεν είναι, ωστόσο, ένα εξωτερικό πλαίσιο που χρησιμοποιείται χωρίς έλεγχο: ο ίδιος ο ορισμός της θέσης ως μονάδας και των οικογενειών της αποτελεί απαιτητική μουσικολογική συνεισφορά, η οποία ελέγχεται σε πραγματικές πηγές και από ανεξάρτητους ειδικούς.

Επειδή η πρόταση απευθύνεται εξίσου σε μουσικολόγους και σε ερευνητές της μηχανικής μάθησης, οι βασικοί υπολογιστικοί όροι ορίζονται εδώ. \emph{Εποπτεία} είναι η πληροφορία από την οποία μαθαίνει ένα μοντέλο. Μπορεί να είναι \emph{ετικέτα}, για παράδειγμα η οικογένεια στην οποία ανήκει μια εμφάνιση, ή \emph{θετικό ζεύγος}, δηλαδή δύο εμφανίσεις που δηλώνονται ως εκδοχές του ίδιου πράγματος· στη δεύτερη περίπτωση μιλάμε για \emph{σχεσιακή εποπτεία}. \emph{Επισημείωση} είναι η καταγραφή αυτών των πληροφοριών από ειδικούς πάνω στις πηγές. Ένας \emph{κωδικοποιητής} μετατρέπει τη συμβολική είσοδο σε \emph{αναπαράσταση}, ένα διάνυσμα αριθμών του οποίου οι αποστάσεις εκφράζουν μαθημένη ομοιότητα. Τα \emph{δεδομένα εκπαίδευσης} εκτιμούν τις παραμέτρους ενός μοντέλου, τα \emph{δεδομένα ανάπτυξης} χρησιμοποιούνται για τις επιλογές του, και τα \emph{δεδομένα ελέγχου} κρίνουν επιλογές που έχουν ήδη γίνει· \emph{μοντέλο αναφοράς} είναι η απλούστερη εναλλακτική που ένα προτεινόμενο μοντέλο οφείλει να ξεπεράσει. Οι μελέτες του κορμού έχουν σύντομους κωδικούς, που ακολουθούν την αρίθμηση του ευρύτερου ερευνητικού προγράμματος: η Μ1 (Μελέτη 1) αφορά την οριοθέτηση και την αναγνώριση, η Α2 είναι το πείραμα για την αξία της μουσικολογικής γνώσης, η Μ2a συγκρίνει τις τρεις συνθήκες B0--B2, η Μ2b εξετάζει συν-εξηγήσεις στην EISE, και η Μ3a παρέχει ανεξάρτητη επικύρωση από ειδικούς· οι Μ3b, Μ3c, η πλήρης Μ4 και η Μ5 περιγράφονται ως διακριτές ερευνητικές επεκτάσεις στο Παράρτημα. Ο Πίνακας~\ref{tab:knowledge} δείχνει πού εισέρχεται η γνώση σε καθεμία.

\begin{table}[htbp]
\caption{Πού εισέρχεται η μουσικολογική γνώση στον κορμό της έρευνας και πώς ελέγχεται}\label{tab:knowledge}
\footnotesize
\begin{tabularx}{\linewidth}{@{}>{\bfseries}L{2.5cm}YYL{1.15cm}@{}}
\toprule
Ρόλος & \textbf{Τι εισάγεται} & \textbf{Πώς ελέγχεται} & \textbf{Μελέτη}\\\midrule
Αναπαράσταση & Μονάδες εισόδου και στρώματα γνωρισμάτων: σημειογραφικά, χρονικά, τροπικά, κειμενικά & Δώδεκα κωδικοποιήσεις με κοινό εκτιμητή, σε έργα εκτός εκπαίδευσης & Μ1\\
Δομικός περιορισμός & Η θέση ως τμήμα μεταβλητού μήκους, με εναλλακτικά όρια, άγνωστες οικογένειες και αποχή & Σύγκριση με γενικές ενδείξεις ορίων και με μαθημένους κατατμητές φράσεων & Μ1\\
Διαγνωστικές σχέσεις & Οι απαιτούμενες, εναλλακτικές και αποκλειόμενες σχέσεις των προδιαγραφών & Έγκυρη ή διαταραγμένη γνώση, σε επίπεδη ή δομημένη χρήση ($2\times2$) & Α2\\
Σχεσιακή εποπτεία & Ζεύγη ανεξάρτητα επισημειωμένων εμφανίσεων της ίδιας οικογένειας & Αναγνώριση και διατήρηση της πραγμάτωσης έναντι κωδικοποιητή χωρίς ζεύγη & Μ2a\\
Ιστορική εποπτεία & Επαληθευμένες συν-εξηγήσεις του ίδιου χωρίου & Έλξη ταυτότητας και σύγκριση τοπικής παραλλαγής με ανεξάρτητο χρονισμό, όπου υπάρχει & Μ2b\\
Εξωτερική επικύρωση & Τυφλές κρίσεις ειδικών για όρια, σχέσεις και ένταξη & Συμφωνία σε νέα χωρία, χωριστά από τα δεδομένα εκπαίδευσης & Μ3a\\\bottomrule
\end{tabularx}
\end{table}

% ============================================================= 2
\section{Μουσικολογική και υπολογιστική θεμελίωση}\label{sec:theory}
Το πρόβλημα της §\ref{sec:intro} βρίσκεται εκεί όπου συναντώνται τρεις ερευνητικές παραδόσεις: η ψαλτική θεωρία και η βυζαντινή μουσικολογία, που όρισαν τη θέση και μελέτησαν την παραλλαγή της· το συμβολικό MIR, που ανέπτυξε μεθόδους για όρια, ομοιότητα και κατηγορία· και η μηχανική μάθηση που αξιοποιεί σχέσεις και προγενέστερη γνώση. Η ενότητα τις παρακολουθεί με αυτή τη σειρά και δείχνει πού συναντώνται.

\subsection{Η θέση στην ψαλτική θεωρία και στη μουσικολογική έρευνα}\label{sec:lit-mus}
Η θέση δεν είναι κατασκευή της σύγχρονης ανάλυσης, γιατί την περιγράφει ήδη η ψαλτική θεωρία. Τον 15ο αιώνα, ο Γαβριήλ ιερομόναχος την περιγράφει με την αναλογία της λέξης (έκδ. Hannick και Wolfram, 1985, σ.~72), ενώ ο Μανουήλ Χρυσάφης ο λαμπαδάριος τη συνδέει με την ένωση σημείων που συγκροτεί μέλος (έκδ. Conomos, 1985, σσ.~40--41). Η αναλογία έχει άμεσο υπολογιστικό περιεχόμενο, αφού τοποθετεί το μουσικό νόημα όχι στο μεμονωμένο σύμβολο αλλά σε ένα ενδιάμεσο επίπεδο, ανάμεσα στο σημείο και στο πλήρες μέλος. Στα χρόνια της μεταρρύθμισης, οι θεωρητικές μαρτυρίες του Αποστόλου Κώνστα του Χίου, σύγχρονου των Τριών Διδασκάλων της Νέας Μεθόδου, του Χρυσάνθου, του Γρηγορίου και του Χουρμουζίου, και εξηγητή ο ίδιος, όπως τις μελέτησε ο Θ.~Κ.~Αποστολόπουλος (2002), δείχνουν πώς συζητούσαν οι ίδιοι οι εξηγητές την απόδοση της παλαιάς γραφής, και ο Χρύσανθος σημειώνει στο \emph{Θεωρητικόν Μέγα} ότι το ίδιο σημάδι αποδίδεται διαφορετικά ανάλογα με τη βαθμίδα, όπως το παρακάλεσμα στους φθόγγους Πα και Βου (1832, §408, κατά την παράθεση του Καραγκούνη, 2003, σ.~715). Η ιστορική θεωρία δίνει έτσι τη μονάδα και τις ονομασίες της, όχι όμως έτοιμες κατηγορίες, και η διαφορά μέλους κάτω από κοινή ονομασία είναι ακριβώς αυτό που πρέπει να ελεγχθεί.

Τον 20ό αιώνα η σχέση της παλαιάς με τη νέα γραφή έγινε συστηματικό ερευνητικό ερώτημα. Ο Κ.~Α.~Ψάχος, στην \emph{Παρασημαντική} του (1917/1978), εξέτασε την ιστορία της σημειογραφίας και το ζήτημα της εξήγησης, ο Γ.~Θ.~Στάθης (1978) μελέτησε συστηματικά τη μετάβαση από την παλαιότερη γραφή στην αναλυτική απόδοση της Νέας Μεθόδου, και ο Σ.~Ι.~Καράς (1982) πρότεινε να διατηρηθούν στη διδασκαλία παλαιά ποιοτικά σημάδια, όπως τα διασώζει η προφορική παράδοση. Από αυτή τη γραμμή έρευνας η διατριβή αντλεί τη διάκριση ανάμεσα στη μετροφωνία και στη μελική ενέργεια των σημαδιών, πάνω στην οποία στηρίζονται οι προδιαγραφές των οικογενειών, καθώς και τους καταγεγραμμένους κανόνες της εξήγησης, που στην ιστορική επέκταση γίνονται μοντέλα αναφοράς.

Η μορφολογική έρευνα των τελευταίων δεκαετιών έδειξε πώς λειτουργεί η θέση μέσα σε συγκεκριμένα ρεπερτόρια. Ο Κ.~Χ.~Καραγκούνης, στη μελέτη του για την παράδοση και την εξήγηση των Χερουβικών και στα \emph{Παραλειπόμενα} (2003, 2005), καταγράφει ονομασμένες θέσεις, επανεμφανίσεις, τροπικές και καταληκτικές προσαρμογές, καθώς και αλυσίδες θέσεων που οργανώνουν τη μορφή, όπως στο Χερουβικό του ψευδο-Ιωάννου Δαμασκηνού, όπου μια αλυσίδα επανέρχεται στο πρώτο τμήμα και η οργάνωση αλλάζει μετά τη λειτουργική τομή της Μεγάλης Εισόδου (2005, σσ.~158--159). Ο Ι.~Αρβανίτης, εξετάζοντας τον ρυθμό μέσα από την παλαιογραφική έρευνα και την εξήγηση (2010), αντιμετωπίζει τη θέση ως αρχή ομοιότητας που αντέχει σε μεταβολές του κειμένου και της γραφής. Από αυτές τις μελέτες προέρχονται τα παραδείγματα της §\ref{sec:unit}, και για τη μάθηση έχουν διπλή αξία, αφού δίνουν υποψήφιες οικογένειες και, κυρίως, τα δύσκολα θετικά και αρνητικά παραδείγματα στα οποία κρίνεται ένα σύστημα.

Τη σημειογραφική και τροπική διάσταση του ερωτήματος τη φωτίζει το έργο της Μ.~Αλεξάνδρου, από την εισαγωγή της στα προβλήματα των εξηγήσεων και των μεταγραφών ως τις λεγόμενες μεγάλες υποστάσεις της μεσοβυζαντινής γραφής και, σε κοινή μελέτη με τον C.~Troelsgård, την τροπικότητα των \emph{Στοιχείων της Παπαδικής} (Αλεξάνδρου, 2010· Alexandru, 2012· Alexandru και Troelsgård, 2020). Μαζί με το εγχειρίδιο του Troelsgård για τα βυζαντινά νεύματα (2011), οι εργασίες αυτές δίνουν το πλαίσιο μέσα στο οποίο ορίζονται οι σημειογραφικές διακρίσεις που χρησιμοποιούν οι προδιαγραφές.

Το πλησιέστερο προηγούμενο για το υλικό της διατριβής είναι η μελέτη του N.~Gheorghiță (2010, σσ.~9--26) για τα κυριακάτικα Κοινωνικά του 18ου αιώνα. Ο Gheorghiță παρακολουθεί έντεκα θέσεις σε 42 Κοινωνικά τριών μελοποιών του 18ου αιώνα, του Δανιήλ Πρωτοψάλτη, του Πέτρου Πελοποννησίου και του Πέτρου Βυζαντίου, και τις αντιπαραβάλλει με τις εξηγήσεις του Γρηγορίου και του Χουρμουζίου. Η εργασία του δείχνει ότι η θέση είναι γόνιμο επίπεδο συγκριτικής ανάλυσης στο παπαδικό ρεπερτόριο, και από αυτήν προέρχονται οι πρώτες οικογένειες του πιλότου. Λίγα χρόνια αργότερα, η διατριβή του Δ.~Σπανουδάκη (2017) για τη μελισματικότητα στα κεκραγάρια, τα μέλη του «Κύριε ἐκέκραξα», συνδύασε παλαιογραφική, δομική και ποσοτική ανάλυση, ενώ στο επίπεδο του πλήρους έργου οι μελέτες για τον Νικόλαο Πρωτοψάλτη Σμύρνης, μελοποιό του 19ου αιώνα (Τόμα, 2018· Ιωαννίδης, 2022), δίνουν ιστορικό πλαίσιο και ένα δημοσιευμένο σημείο ελέγχου για μελλοντική εφαρμογή. Η διατριβή μεταφέρει αυτή τη μουσικολογική μέθοδο σε αναπαραγώγιμο υπολογιστικό πρωτόκολλο, με προδιαγραφές που εφαρμόζονται σε ολόκληρο σώμα και με ανεξάρτητα κριτήρια.

\subsection{Από την κωδικοποίηση των νευμάτων στην οριοθέτηση και την ταυτότητα}\label{sec:lit-mir}
Η υπολογιστική έρευνα της βυζαντινής ψαλτικής ξεκίνησε νωρίς, αλλά προχώρησε με λίγα και ουσιαστικά βήματα. Ήδη ο N.~Schiødt (1971) ανέλυσε με υπολογιστή τριάντα πέντε βυζαντινούς ύμνους, και ο P.~Mavromatis διατύπωσε αργότερα πιθανοτικό μοντέλο της μελωδικής παραγωγής στη σύγχρονη ελληνική ψαλτική με κρυφά μαρκοβιανά μοντέλα (2005a), εξετάζοντας στη διατριβή του τις γνωσιακές του συνέπειες για τη γραπτή και την προφορική μετάδοση (2005b). Η τυποποίηση της ίδιας της γραφής προχώρησε με την κωδικοποίηση της βυζαντινής σημειογραφίας (Bouris κ.ά., 2018), με τα συγκρίσιμα ψηφιακά σώματα μεσοβυζαντινής γραφής των P.~Polykarpidis κ.ά. (2022) και με την ιεραρχική αναπαράσταση της ίδιας γραφής στη διατριβή του Π.~Πολυκαρπίδη για το Ειρμολόγιο του Θεοφάνη Καρύκη (2023). Πρόσφατα, η Αλεξάνδρου, με τους E.~De Luca και I.~Fujinaga, εξέτασε τη διαλειτουργική κωδικοποίηση των μεσοβυζαντινών νευμάτων στο πρότυπο του Music Encoding Initiative (MEI), με προσοχή στο πώς η σημασία ενός νεύματος εξαρτάται από το συμφραζόμενό του (Alexandru κ.ά., 2024), και μια πιλοτική κωδικοποίηση 32 νευμάτων της γραφής της Νέας Μεθόδου παρουσιάστηκε στο συνέδριο της Διεθνούς Ένωσης Μουσικών Βιβλιοθηκών, Αρχείων και Κέντρων Τεκμηρίωσης (IAML· Thomae κ.ά., 2026). Οι εργασίες αυτές προωθούν την αναπαράσταση στο επίπεδο του σημείου, ενώ η οριοθέτηση και η αναγνώριση θέσεων μεταβλητής έκτασης μέσα στη συνεχή γραφή του παπαδικού μέλους δεν έχει, εξ όσων εντοπίσαμε, αντιμετωπιστεί.

Η μελωδική κατάτμηση, δηλαδή ο αυτόματος εντοπισμός ορίων μέσα σε μια μελωδία, έχει μακρά ιστορία στο συμβολικό MIR, και η εξέλιξή της δείχνει καθαρά τι λείπει. Τα πρώτα μοντέλα τοποθετούσαν το όριο σε τοπικές ασυνέχειες τονικού ύψους και διάρκειας, όπως το Local Boundary Detection Model (LBDM· Cambouropoulos, 2001a), ή στα άκρα επαναλαμβανόμενων μελωδικών δομών (Cambouropoulos, 2006). Τα πληροφοριακά μοντέλα, με κύριο εκπρόσωπο το IDyOM (Information Dynamics of Music), που εκτιμά πόσο απρόσμενο είναι κάθε επόμενο γεγονός, έδειξαν αργότερα ότι η μουσική προσδοκία προβλέπει τα όρια που αντιλαμβάνονται οι ακροατές (Pearce, Müllensiefen και Wiggins, 2010), και την τελευταία δεκαετία τα όρια φράσεων μαθαίνονται από επισημειωμένα δεδομένα με συνελικτικά και αναδρομικά νευρωνικά δίκτυα που καταλήγουν σε τυχαίο πεδίο υπό συνθήκη (conditional random field, CRF), ένα μοντέλο που αποφασίζει μαζί τις ετικέτες όλης της ακολουθίας (Guan κ.ά., 2018), με εξόρυξη κανόνων και τυχαία δάση (van~Kranenburg, 2020) και με Transformer, νευρωνικά δίκτυα που συσχετίζουν κάθε γεγονός με όλα τα υπόλοιπα μέσω μηχανισμού προσοχής (Guan, Dong, Liu και Li, 2025). Στα αντίστοιχα σώματα αξιολόγησης τα μαθημένα μοντέλα παρουσιάζουν βελτιώσεις, ενώ το LBDM και το IDyOM παραμένουν τα τυπικά ερμηνεύσιμα μοντέλα αναφοράς (van~Kranenburg, 2020· Guan κ.ά., 2025). Οι συγκεκριμένες εργασίες αξιολογούν κυρίως όρια φράσεων και προσφέρουν συγκριτές για τη χωριστή συνιστώσα της οριοθέτησης. Στο γρηγοριανό μέλος, η μη εποπτευόμενη κατάτμηση εξέτασε ποιες μονάδες αναδύονται από τη στατιστική του σώματος (Cornelissen κ.ά., 2020· Lanz και Hajič, 2025), και οι Lanz και Hajič έδειξαν ότι τα τμήματα αυτά δεν συμπίπτουν αναγκαστικά με τις μονάδες που αναγνωρίζει η μουσικολογία. Καθώς και οι ίδιοι οι επισημειωτές διαφωνούν συχνά για τα όρια (Wang κ.ά., 2017), η μονάδα χρειάζεται εξωτερικό ορισμό και περισσότερα από ένα αποδεκτά όρια. Τα ημι-μαρκοβιανά CRF (semi-Markov CRF, semi-CRF), που βαθμολογούν ολόκληρα τμήματα μαζί με την ετικέτα τους (Sarawagi και Cohen, 2004), εκπαιδεύονται από μερικές επισημειώσεις, όπου δεν είναι γνωστά όλα τα όρια (Tsuboi κ.ά., 2008) και δέχονται νευρωνικές αναπαραστάσεις τμημάτων (Kong, Dyer και Smith, 2016), προσφέρουν το φυσικό εργαλείο για να εκτιμώνται μαζί η έκταση και η κατηγορία.

Η ταυτότητα αποδείχθηκε δυσκολότερη από το όριο. Η ομοιότητα και ο σχηματισμός κατηγοριών αλληλο\-εξαρτώνται, αφού η βαρύτητα κάθε γνωρίσματος εξαρτάται από τις κατηγορίες που συγκροτούνται (Cambouropoulos, 2001b). Στις οικογένειες μελωδιών των ολλανδικών παραδοσιακών τραγουδιών, οι A.~Volk και P.~van~Kranenburg (2012) έδειξαν ότι οι ειδικοί κατατάσσουν με ρητά κριτήρια, αλλά κανένα γνώρισμα δεν λειτουργεί ως καθολικός πυρήνας για όλες τις οικογένειες, ενώ ο A.~Holzapfel (2010) παρατήρησε ότι η ομοιότητα δύο αποσπασμάτων οφείλεται συχνά σε ένα μικρό κοινό υποτμήμα και όχι σε ταυτότητα ολόκληρης της μονάδας. Πάνω σε αυτή τη γνώση στηρίχθηκαν μέθοδοι ανακάλυψης τμηματικών προτύπων (Conklin και Anagnostopoulou, 2006), εντοπισμού μελωδικών τμημάτων (Janssen κ.ά., 2017) και μάθησης μετρικών συγγένειας (Karsdorp κ.ά., 2019). Το πιο κοντινό ανάλογο στη δυτική μουσική είναι η αναγνώριση των σχημάτων του γκαλάν ύφους, κατηγοριών που ορίζει η θεωρία και όχι η στατιστική (Gjerdingen, 2007· Finkensiep κ.ά., 2020), και οι G.~Rabinovitch και A.~Carter-Enyi (2024) έδειξαν ότι η εσωτερική σταθερότητα διαφέρει από σχήμα σε σχήμα, αξιοποιώντας μάλιστα τα μοντέλα του Mavromatis για τους τύπους της ελληνικής ψαλτικής. Αυτό που λείπει είναι η σύνδεση αυτών των γραμμών: η αναγνώριση μονάδων με εξωτερικά ορισμένη ταυτότητα, μεταβλητό και αμφίσημο μήκος και ρητές, ελέγξιμες μουσικολογικές προδιαγραφές.

Η συγκριτική υπολογιστική έρευνα άλλων τροπικών παραδόσεων δείχνει ότι η μελωδική μονάδα μπορεί να οριστεί και να αναζητηθεί πέρα από το επίπεδο της κλίμακας ή του τρόπου. Στην ινδική μουσική \emph{rāga}, οι χαρακτηριστικές φράσεις \emph{pakad} έχουν χρησιμοποιηθεί για την ταξινόμηση μουσικών μοτίβων και τη μελέτη των διαφορετικών πραγματώσεών τους σε ηχογραφήσεις, με δεδομένα ήδη χωρισμένα σε φράσεις (Rao κ.ά., 2014). Στην Arab-Andalusian μουσική, η ανακάλυψη μελωδικών προτύπων συσχετίστηκε με τα centos που αναγνωρίζει ειδικός της παράδοσης (Nuttall κ.ά., 2021), ένα κοντινό προηγούμενο πολιτισμικά τεκμηριωμένης μονάδας. Στην οθωμανική/τουρκική παράδοση \emph{makam}, η αυτόματη κατάτμηση έχει συνδυάσει γενικές μελωδικές ενδείξεις με γνωρίσματα της τροπικής και ρυθμικής οργάνωσης (Bozkurt κ.ά., 2014). Τα αποτελέσματα αυτά, μαζί με τη συμβολική κατάτμηση του γρηγοριανού μέλους που προαναφέρθηκε, αποτελούν ουσιώδη προηγούμενα και όχι ισοδύναμα της βυζαντινής θέσης: η \emph{pakad}, η φράση ενός \emph{makam} και η θέση έχουν διαφορετικό θεωρητικό και λειτουργικό καθεστώς. Η παρούσα έρευνα αξιοποιεί αυτή την εμπειρία για να ελέγξει ένα συγκεκριμένο, σύνθετο ερώτημα: αν μια ιστορικά προσδιορισμένη μουσική μονάδα μπορεί να οριοθετηθεί και να αναγνωριστεί μέσα σε συνεχή συμβολική γραφή βάσει ρητών, ανεξάρτητα ελεγχόμενων κριτηρίων ταυτότητας, και τι προσθέτει η γνώση αυτή στη μάθηση.

\subsection{Μάθηση από σχέσεις και από δομημένη γνώση}\label{sec:lit-ml}
Για να μάθει ένα μοντέλο αυτή την ταυτότητα, πρέπει να γνωρίζει ποια παραδείγματα είναι εκδοχές του ίδιου πράγματος, και η μηχανική μάθηση έχει δείξει ότι η επιλογή αυτή καθορίζει το αποτέλεσμα. Στην εποπτευόμενη αντιπαραθετική μάθηση, που φέρνει κοντά στην αναπαράσταση τα παραδείγματα που δηλώνονται ως ίδια και απομακρύνει τα υπόλοιπα, τα θετικά ζεύγη προέρχονται από κοινές ετικέτες κατηγορίας (Khosla κ.ά., 2020), και η προέλευσή τους ορίζει ποιες διαφορές θα αγνοήσει η αναπαράσταση· η αυτοεπιβλεπόμενη μάθηση, από την άλλη, όπου το μοντέλο εκπαιδεύεται να συμπληρώνει κρυμμένα γεγονότα, μαθαίνει από τη συμβολική μουσική χωρίς ετικέτες (Zeng κ.ά., 2021). Για ένα μουσικολογικό πρόβλημα το κρίσιμο ερώτημα είναι επομένως διπλό: τι κερδίζει η αναγνώριση από τη σχέση κοινής οικογένειας, και τι χάνεται από την πληροφορία της συγκεκριμένης πραγμάτωσης.

Η μουσικολογική γνώση δεν δίνει όμως μόνο ζεύγη· δίνει και κανόνες με λογική δομή, με απαιτούμενες σχέσεις, εναλλακτικούς δρόμους και αποκλεισμούς. Όταν η απόφαση περνά από ενδιάμεσες έννοιες με σαφές νόημα, όπως στα μοντέλα στενωπού εννοιών (concept bottleneck models· Koh κ.ά., 2020), γίνεται ερμηνεύσιμη, και η σημασιολογική απώλεια (semantic loss), ο αρνητικός λογάριθμος της πιθανότητας των αναθέσεων που ικανοποιούν έναν κανόνα, επιτρέπει να επιβάλλεται ένας λογικός κανόνας κατά τη μάθηση (Xu κ.ά., 2018). Η νευροσυμβολική βιβλιογραφία, που συνδυάζει νευρωνικά δίκτυα με λογικούς κανόνες, έχει αναδείξει όμως δύο κινδύνους: ένα μοντέλο μπορεί να ικανοποιεί τον κανόνα μαθαίνοντας λανθασμένες ενδιάμεσες έννοιες (Marconato κ.ά., 2023), και η υπόθεση ότι οι έννοιες είναι ανεξάρτητες παραμορφώνει την πιθανότητα (van~Krieken κ.ά., 2024). Και οι δύο κίνδυνοι οδηγούν στο ίδιο μεθοδολογικό συμπέρασμα, που διατρέχει τη διατριβή: η βελτίωση από την ενσωμάτωση της γνώσης πρέπει να συγκρίνεται με τη βελτίωση από μια ελεγχόμενα λανθασμένη γνώση ίδιας μορφής, αλλιώς δεν είναι δυνατόν να ξέρουμε τι ακριβώς ωφέλησε. Η ίδια λογική ισχύει για την αξιολόγηση: η βιβλιογραφία για τη διαρροή, δηλαδή για πληροφορία από τα δεδομένα ελέγχου που επηρεάζει ακούσια την ανάπτυξη ενός μοντέλου, δείχνει πόσο εύκολα αλλοιώνεται ένα αποτέλεσμα όταν επιλογές μοντελοποίησης γίνονται με δεδομένα που αργότερα χρησιμοποιούνται για έλεγχο (Kapoor και Narayanan, 2023), και στο MIR ο Sturm (2014) έδειξε ότι υψηλή επίδοση μπορεί να στηρίζεται σε γνωρίσματα άσχετα με τη μουσική έννοια που υποτίθεται ότι μετριέται, δηλαδή σε \emph{συντομεύσεις}.

% ============================================================= 3
\section{Ερευνητικό κενό, ερωτήματα και συνεισφορά}\label{sec:rq}
Η επισκόπηση δείχνει ότι τα συστατικά του προβλήματος υπάρχουν, αλλά σε χωριστές ερευνητικές γραμμές. Η ιστορικά τεκμηριωμένη συνεξήγηση προσθέτει διαφορετικό είδος εποπτείας: δύο αποδόσεις συνδέονται επειδή αφορούν το ίδιο τεκμηριωμένο χωρίο και όχι επειδή ο αλγόριθμος τις έκρινε απλώς όμοιες. Η μελωδική κατάτμηση έχει ώριμες μεθόδους για όρια φράσεων χωρίς ταυτότητα, η έρευνα των οικογενειών μελωδιών έχει μεθόδους ομοιότητας χωρίς οριοθέτηση, η νευροσυμβολική μάθηση ενσωματώνει κανόνες χωρίς να ελέγχει συνήθως το περιεχόμενό τους, και η βυζαντινή μουσικολογία έχει ορίσει τη θέση και τις δύσκολες περιπτώσεις της με ποιοτικές συγκρίσεις έργο προς έργο. Παρά τα σημαντικά προηγούμενα σε τροπικές μουσικές παραδόσεις, δεν εντοπίσαμε μελέτη που να συνδυάζει, στο επίπεδο της βυζαντινής θέσης, κοινή οριοθέτηση και αναγνώριση μεταβλητών τμημάτων, ανεξάρτητα ελεγχόμενες οικογενειακές προδιαγραφές, εναλλακτά αποδεκτά όρια, άγνωστες οικογένειες και πειραματική αποτίμηση του περιεχομένου της μουσικολογικής γνώσης. Η θέση είναι κατάλληλο πεδίο ακριβώς επειδή όλα αυτά είναι εδώ τεκμηριωμένα. Το κεντρικό ερώτημα της διατριβής διατυπώνεται έτσι:

\begin{quote}\itshape
Πώς αναγνωρίζεται μέσα στη συνεχή γραφή μια μουσική θέση και η οικογένειά της όταν η πραγμάτωση μεταβάλλεται, και τι προσφέρει σε αυτή την αναγνώριση η μουσικολογική γνώση που ορίζει την οικογένεια;
\end{quote}

\noindent Το ερώτημα αναλύεται σε τρία ερωτήματα, που χρησιμοποιούν τα ίδια δεδομένα, τους ίδιους διαχωρισμούς και την ίδια αναπαράσταση εισόδου.

\begin{description}[leftmargin=0pt,itemsep=1pt,parsep=0pt,font=\normalfont\bfseries]
\item[ΕΕ1 --- Κοινή οριοθέτηση και αναγνώριση (Μ1).] Ποια μονάδα εισόδου και ποια στρώματα γνωρισμάτων υποστηρίζουν καλύτερα την κοινή εκτίμηση ορίων και οικογένειας, και πόση πληροφορία προσθέτουν σε αυτήν οι γενικές ενδείξεις μελωδικής κατάτμησης;
\item[ΕΕ2 --- Περιεχόμενο και μορφή της γνώσης (Α2).] Βελτιώνει η έγκυρη προδιαγραφή την αναγνώριση της οικογένειας σε σχέση με μια ελεγχόμενα διαταραγμένη προδιαγραφή ίδιας μορφής, και τι προσθέτει η λογική σύνθεση των έγκυρων διαγνωστικών σχέσεων στην προσθετική χρήση τους;
\item[ΕΕ3 --- Ταυτότητα, πραγμάτωση και ιστορική εποπτεία (Μ2).] Τι προσφέρει η μάθηση από εμφανίσεις ίδιας οικογένειας στην αναγνώριση και στη διατήρηση της παραλλαγής (B0--B2); Στο επιλέξιμο ιστορικό πεδίο, πώς αλλάζει η συμφωνία της τοπικής παραλλαγής με ανεξάρτητο εκτελεστικό χρονισμό όταν τα ζεύγη είναι συν-εξηγήσεις του ίδιου χωρίου (B2--EISE);
\end{description}

Από τα ερωτήματα προκύπτουν τρεις συνεισφορές. Η πρώτη (\textbf{C1}) είναι μια νέα εργασία συμβολικού MIR, η κοινή οριοθέτηση και αναγνώριση οικογενειών θέσεων με αμφίσημα όρια και άγνωστες οικογένειες, μαζί με πρωτόκολλο και σώμα αξιολόγησης. Η δεύτερη (\textbf{C2}) είναι ένας πειραματικός σχεδιασμός που μετρά χωριστά την αξία του περιεχομένου και της μορφής της μουσικολογικής γνώσης στη μάθηση. Η τρίτη (\textbf{C3}) εξετάζει δύο συμπληρωματικές μορφές μάθησης αναπαραστάσεων: τη σχεσιακή μάθηση οικογένειας με διατήρηση της πραγμάτωσης (\textbf{C3a}, Μ2a) και την ιστορικά τεκμηριωμένη εποπτεία από συν-εξηγήσεις (\textbf{C3b}, Μ2b), με χωριστή διαγνωστική και, εφόσον τα δεδομένα επαρκούν, ανεξάρτητη εκτελεστική αξιολόγηση. Η μουσικολογική συμβολή δεν εξαντλείται στο συνοδευτικό σώμα επισημειώσεων: περιλαμβάνει τον ελέγξιμο ορισμό της θέσης ως μονάδας, την αιτιολόγηση των οικογενειακών διακρίσεων και την τεκμηριωμένη καταγραφή της παραλλαγής και της αμφισημίας τους. Ο Πίνακας~\ref{tab:questions} συνδέει κάθε ερώτημα με το κύριο μέγεθος, τη σύγκριση και το καθεστώς του συμπεράσματος. Τα σύμβολά του ορίζονται αναλυτικά στις §§\ref{sec:m1}--\ref{sec:repr}· συνοπτικά, $U\times\Phi$ είναι οι δώδεκα αναγνώσεις της γραφής που συγκρίνει η Μ1 και $\Rs$ εκείνη που επιλέγεται, $F1_{\mathrm{seg}}$ μετρά πόσα τμήματα βρίσκονται με σωστά όρια και σωστή οικογένεια, $F1_{\mathrm{bnd}}$ μόνο τα όρια και $\mathrm{Rec}_{\mathrm{fam}}$ τη μέση ανάκληση ανά οικογένεια, $\mathrm{AP}_{\mathrm{fam}}$ μετρά την κατάταξη των εμφανίσεων ανά οικογένεια, με γνωστά όρια, και κατόπιν λαμβάνεται ο μέσος όρος των οικογενειών, τα $\Delta_{\mathrm{know}}$, $\Delta_{\mathrm{logic},V}$ και $\Delta_{\mathrm{famrel}}$ είναι διαφορές αυτής της μέτρησης ανάμεσα σε συνθήκες, ενώ το $\dhist$ είναι διαφορά μέσων συσχετίσεων κατάταξης, το $\Qreal$ μετρά πόση πληροφορία της συγκεκριμένης πραγμάτωσης διατηρείται, και τα $G$ με δείκτη είναι τα κριτήρια επάρκειας των δεδομένων (§\ref{sec:gates}). \emph{Επιβεβαιωτικό} καθεστώς σημαίνει ότι η υπόθεση και το κριτήριο δηλώνονται πριν από τα αποτελέσματα και ελέγχονται στατιστικά· \emph{περιγραφικό}, ότι τα μεγέθη αναφέρονται με την αβεβαιότητά τους χωρίς επιβεβαιωτικό ισχυρισμό. Το Σχήμα~\ref{fig:core} δείχνει πώς συνδέονται οι μελέτες.

\begin{table}[htbp]
\caption{Ερωτήματα, κύρια μεγέθη, συγκρίσεις και καθεστώς}\label{tab:questions}
\footnotesize
\begin{tabularx}{\linewidth}{@{}L{0.85cm}L{3.3cm}YL{2.7cm}@{}}
\toprule
\textbf{ΕΕ} & \textbf{Κύριο μέγεθος} & \textbf{Σύγκριση και μοντέλα αναφοράς} & \textbf{Καθεστώς}\\\midrule
ΕΕ1 & $F1_{\mathrm{seg}}$ τμήματος με σωστή οικογένεια· $F1_{\mathrm{bnd}}$ ορίων· $\mathrm{Rec}_{\mathrm{fam}}$ & Δώδεκα αναπαραστάσεις σε A/B/C· μικρός Transformer, $T_{\mathrm{evt}}$, Pitman--Yor και μουσικές ενδείξεις ορίων & Προκαθορισμένο περιγραφικό· επιλογή $\Rs$ με κανόνα\\
ΕΕ2 & $\Delta_{\mathrm{know}}$, $\Delta_{\mathrm{logic},V}$ στο $\mathrm{AP}_{\mathrm{fam}}$ & $2\times2$: έγκυρη ή διαταραγμένη προδιαγραφή $\times$ επίπεδη ή δομημένη χρήση, στα ίδια όρια και την ίδια $\Rs$ & Επιβεβαιωτικό με τεχνική και στατιστική επάρκεια· αλλιώς εκτιμητικό\\
ΕΕ3a & $\Delta_{\mathrm{famrel}}$· $\Qreal$ & Μ2a: B2 έναντι B1· B0 συμβολική αναφορά & Περιγραφικό, έλεγχος διατήρησης\\\midrule
ΕΕ3b & $\dhist$· $S_{\mathrm{EISE}},S_{\mathrm{B2}}$ & Μ2b: B2--EISE· ανεξάρτητο $\Gamma_f(u)$ (Μ4-min) & Επιβεβαιωτικό μόνο με τα κριτήρια επάρκειας\\\bottomrule
\end{tabularx}
\end{table}

\begin{figure}[htbp]
\centering
\begin{tikzpicture}[font=\footnotesize]
\node[boxa,text width=3.7cm,minimum height=1.1cm] (src) at (0,0) {\textbf{ΠΕ0 και πιλότος}\\πηγές, προδιαγραφές, επισημείωση, κριτήρια επάρκειας};
\node[boxa,text width=3.8cm,minimum height=1.1cm] (m1) at (4.75,0) {\textbf{Μ1}\\όρια και οικογένεια μαζί·\\επιλογή $\Rs$};
\node[boxa,text width=3.4cm,minimum height=1.1cm] (a2) at (9.0,0.75) {\textbf{Α2}\\έγκυρη ή διαταραγμένη γνώση $\times$ χρήση};
\node[boxa,text width=3.4cm,minimum height=1.1cm] (b) at (9.0,-0.75) {\textbf{Μ2a: B0--B2}\\ταυτότητα και πραγμάτωση};
\node[boxb,text width=3.0cm,minimum height=1.1cm] (h) at (13.3,-0.75) {\textbf{Μ2b: B2--EISE}\\ιστορικά ζεύγη· $\dhist$ υπό κριτήρια};
\node[boxa,text width=3.7cm,minimum height=0.8cm] (m3) at (0,-1.6) {\textbf{Μ3a}\\τυφλή επικύρωση};
\draw[arr] (src) -- (m1);
\draw[arr] (m1.east) -- ++(0.35,0) |- (a2.west);
\draw[arr] (m1.east) -- ++(0.35,0) |- (b.west);
\draw[arr] (b) -- (h);
\draw[arr] (src) -- (m3);
\draw[arr,dashed] (m3.east) -| node[pos=0.3,above,font=\scriptsize]{κρίσεις εκτός εκπαίδευσης} (m1.south);
\end{tikzpicture}
\caption{Η κοινή αρχιτεκτονική της έρευνας. Η Μ1 επιλέγει την αναπαράσταση $\Rs$, ενώ η Α2 και η Μ2a εξετάζουν ανεξάρτητα την αξία της μουσικολογικής γνώσης και τη μάθηση οικογενειακής ταυτότητας. Η Μ2b συνεχίζει τη Μ2a με ιστορικά τεκμηριωμένη εποπτεία: η απογραφή συν-εξηγήσεων είναι προγραμματισμένο παραδοτέο, ενώ η EISE εκπαιδεύεται με επιλέξιμα ζεύγη και το $\dhist$ ελέγχεται μόνο με επαρκή ανεξάρτητα εκτελεστικά δεδομένα (Μ4-min) και στατιστική ακρίβεια. Η Μ3a παρέχει ανεξάρτητες μουσικολογικές κρίσεις σε όλη τη μελέτη.}
\label{fig:core}
\end{figure}

% ============================================================= 4
\section{Σώμα δεδομένων, πιλότος και κριτήρια επάρκειας}\label{sec:data}
\subsection{Πηγές, διαθεσιμότητα και δικαιώματα}\label{sec:corpus}
Και τα τρία ερωτήματα προϋποθέτουν ένα σώμα πλήρων έργων σε συμβολική μορφή, με επισημειωμένες εμφανίσεις θέσεων. Το μουσικό πεδίο είναι τα παπαδικά Χερουβικά και Κοινωνικά της Νέας Μεθόδου. Ο δημόσιος κατάλογος \emph{Papadic Formulae} της Ιεράς Μονής Αγίου Αντωνίου στην Αριζόνα συγκεντρώνει τύπους από 758 συνθέσεις (349 Χερουβικά, 351 Κοινωνικά, 58 λοιπά) και χρησιμεύει ως οδηγός για τον εντοπισμό υλικού και για υποψήφιες οικογένειες· δεν είναι όμως έτοιμο, μηχανικά αναγνώσιμο σώμα πλήρων έργων. Τα έργα κωδικοποιούνται από αρχεία του Neanes, ανοικτού λογισμικού γραφής της βυζαντινής σημειογραφίας (Neanes Project, χ.χ.), και ο αριθμός των Χερουβικών και Κοινωνικών που είναι ήδη διαθέσιμα σε πλήρη, επαληθεύσιμη μορφή παραμένει ανοικτό στοιχείο της απογραφής. Το προπαρασκευαστικό πακέτο εργασίας ΠΕ0 (μήνες 1--6) θα τον καταγράψει ανά έργο, ήχο, μελοποιό, πηγή και καθεστώς άδειας πριν επιλεγεί το τελικό δείγμα.

Κάθε πηγή μετατρέπεται σε ενιαίο συμβολικό επίπεδο αναφοράς, μια ακολουθία ατομικών μελωδικών γεγονότων με αμφίδρομη σύνδεση προς τα γραφικά σημεία, ώστε κάθε αυτόματη απόφαση να ελέγχεται στο χωρίο που την παρήγαγε. Υπάρχει ήδη ερευνητική υποδομή μετατροπής, η οποία θα αξιολογηθεί σε πραγματικά έργα με έλεγχο αντίστροφης ανακατασκευής, δηλαδή επαναφοράς της αρχικής γραφής από τη συμβολική μορφή, και με καταγραφή των πεδίων που χάνονται. Η άδεια του λογισμικού διακρίνεται από τα δικαιώματα των μουσικών αρχείων, και γι' αυτό το ΠΕ0 απογράφει χωριστά τα δικαιώματα επεξεργασίας, δημοσίευσης μεταγραφών και επαναδιανομής επισημειώσεων. Το σώμα αξιολόγησης θα διατεθεί δημόσια όπου το επιτρέπουν οι άδειες· διαφορετικά δημοσιεύονται οι επισημειώσεις, τα ευρετήρια και το πρωτόκολλο, με παραπομπή στις πηγές.

\subsection{Πιλότος μηνών 1--6: όλη η αλυσίδα σε μικρή κλίμακα}\label{sec:pilot}
Πριν από οποιαδήποτε δέσμευση σε μεγάλη κλίμακα, ο πιλότος διατρέχει όλη την αλυσίδα σε τρεις υποψήφιες οικογένειες, την παρακλητική, το κύλισμα και το πίασμα, με αφετηρία τις συγκρίσεις του Gheorghiță (2010) σε Κοινωνικά. Περιλαμβάνει επιλογή πηγών, μετατροπή και έλεγχο ανακατασκευής, δύο ανεξάρτητες επισημειώσεις ορίων και ένταξης, ένα βασικό semi-CRF και αξιολόγηση σε χωριστά έργα. Περιλαμβάνει επίσης την πρώτη δοκιμή της Α2: μία προδιαγραφή γράφεται σε λογική μορφή (§\ref{sec:a2}) και οι έννοιές της υπολογίζονται αυτόματα από τη συμβολική είσοδο, ώστε να φανεί πόσες από αυτές μπορούν να εξαχθούν αξιόπιστα. Καταγράφονται ο χρόνος ανά έργο και ανά εμφάνιση, το ποσοστό έγκυρης μετατροπής, η συμφωνία των επισημειωτών και η συχνότητα των δύσκολων περιπτώσεων. Παράλληλα απογράφονται οι ανεξάρτητες ιστορικές γραμμές συν-εξηγήσεων και η πραγματική τομή τους με ηχογραφήσεις· χωρίς επαρκές υλικό η καταγραφή παραμένει αυτοτελές παραδοτέο, όχι υπόσχεση επιβεβαιωτικού πειράματος.

Οι αριθμοί που ακολουθούν είναι στόχοι σχεδιασμού προς επιβεβαίωση από τον πιλότο και όχι καταμετρημένη διαθεσιμότητα. Το τελικό σώμα στοχεύει σε 8--12 οικογένειες, με τουλάχιστον 25--30 εμφανίσεις ανά οικογένεια από τουλάχιστον πέντε ανεξάρτητα έργα, σε 40--60 έργα συνολικά. Τουλάχιστον το 30\% των εμφανίσεων και όλες οι δύσκολες περιπτώσεις επισημειώνονται τυφλά από δεύτερο ειδικό, με στόχο συμφωνίας ένταξης, μετρημένης με τον συντελεστή κ του Cohen, $\kappa\geq0{,}70$. Ο προϋπολογισμός εργασίας χωρίζεται σε τέσσερις γραμμές, για τις οποίες η αρχική υπόθεση εργασίας είναι η εξής: έλεγχος μεταγραφής και πρώτη επισημείωση 2--4 ώρες ανά πλήρες έργο, δεύτερη τυφλή επισημείωση περίπου 1--2 ώρες ανά έργο του υποσυνόλου, επίλυση διαφωνιών, και επικύρωση σε δείγμα των εξαγωγέων εννοιών, δηλαδή των προγραμμάτων που υπολογίζουν τις διαγνωστικές έννοιες. Με 40--60 έργα, η πρώτη γραμμή αντιστοιχεί σε περίπου 80--240 ώρες. Οι εκτιμήσεις αντικαθίστανται στον μήνα 6 από τις μετρήσεις του πιλότου, και ο ρόλος, ο χρόνος και η αναγνώριση της συμβολής των δύο ειδικών συμφωνούνται ρητά στο ΠΕ0.

\subsection{Κριτήρια επάρκειας και προκαθορισμένες αποφάσεις}\label{sec:gates}
Στον έκτο μήνα τέσσερα κριτήρια επάρκειας, που συμβολίζονται με $G$ και δείκτη, καθορίζουν την κλίμακα του κορμού (Πίνακας~\ref{tab:gates}). Η ιστορική Μ2b διαθέτει επιπλέον τα $G_{\mathrm{pair}}$, για τεκμηριωμένα και μη εξαρτημένα ζεύγη συν-εξηγήσεων, $G_{\mathrm{perf}}$, για ηχογραφήσεις των ίδιων χωρίων από ανεξάρτητους ψάλτες, και, μετά την επιλογή $\Rs$, το $G_{\mathrm{precision}}$, για βαθμονομημένη ακρίβεια της σύγκρισης· χωρίς αυτά δεν ενεργοποιείται η επιβεβαιωτική ιστορική αντίθεση. Τα κριτήρια σχηματίζουν μια κλιμάκωση: ο στόχος είναι 8--12 οικογένειες· το ελάχιστο για τη Μ1 και τις B0--B2 είναι έξι· και η τεχνική Α2 απαιτεί τουλάχιστον τέσσερις οικογένειες με εξαγώγιμες έννοιες και ενεργές, μη ισοδύναμες διαταραχές. Η στατιστική επιβεβαίωση απαιτεί επιπλέον επαρκείς ανεξάρτητες ομάδες και ελεγμένη συμπεριφορά της συμπερασματολογίας· ο αριθμός οικογενειών μόνος του δεν αρκεί. Στον ίδιο μήνα ελέγχεται η εφικτότητα των ιστορικών $G_{\mathrm{pair}}$ και $G_{\mathrm{perf}}$. Η τελική στατιστική ενεργοποίηση της §\ref{sec:historicpriority} επανελέγχεται μετά την επιλογή $\Rs$ και πριν από το ιστορικό σύνολο ελέγχου· η απογραφή και η τεκμηρίωση σχέσεων παραμένουν παραδοτέα ανεξάρτητα από αυτήν.

\begin{table}[htbp]
\caption{Κριτήρια επάρκειας του κορμού: στόχοι σχεδιασμού και αποφάσεις σε περίπτωση ανεπάρκειας}\label{tab:gates}
\footnotesize
\begin{tabularx}{\linewidth}{@{}L{1.3cm}YL{4.6cm}@{}}
\toprule
\textbf{Κριτήριο} & \textbf{Αρχικός στόχος, προς επικύρωση στο ΠΕ0} & \textbf{Απόφαση όταν το κριτήριο δεν ικανοποιείται}\\\midrule
$G_{\mathrm{symbol}}$ & Τουλάχιστον 40 πλήρη έργα με ελέγξιμη μετατροπή και σαφές δικαίωμα επεξεργασίας & Περιορισμός σε Κοινωνικά ή σε επιλεγμένους ήχους· στοχευμένη πρόσθετη κωδικοποίηση· νέα απογραφή\\
$G_{\mathrm{fam}}$ & Τουλάχιστον 6 οικογένειες με 25 εμφανίσεις από $\geq5$ ανεξάρτητα έργα και συμφωνία $\kappa\geq0{,}70$ & Περιορισμός στις σταθερότερες οικογένειες· συγχώνευση μόνο με μουσικολογική αιτιολόγηση· οι υπόλοιπες μένουν διερευνητικές\\
$G_{\mathrm{split}}$ & Αρκετές ανεξάρτητες ομάδες έργων ώστε κάθε οικογένεια να εμφανίζεται σε κάθε γύρο της διασταυρούμενης επικύρωσης & Λιγότεροι γύροι ή αποκλεισμός μίας ομάδας έργων κάθε φορά· αναφορά της αυξημένης αβεβαιότητας\\
$G_{\mathrm{know}}$ & Τουλάχιστον 4 προδιαγραφές με εξαγώγιμες έννοιες, 10 ενεργές διαταραχές με πιστοποιητικό μη ισοδυναμίας· χωριστός έλεγχος στατιστικής επάρκειας & Η Α2 εκτελείται ως διερευνητική σύγκριση στις διαθέσιμες οικογένειες· η Μ1 και οι B0--B2 συνεχίζουν αυτοτελώς\\\bottomrule
\end{tabularx}
\end{table}

\subsection{Σχεδιασμός αξιολόγησης και ρόλοι των τεκμηρίων}\label{sec:design}
Οι ομάδες ορίζονται από συγγένεια έργων και πηγών, ώστε αντίγραφα και συν-εξηγήσεις να παραμένουν μαζί. Η ομαδοποιημένη εμφωλευμένη διασταυρούμενη επικύρωση εκτιμά ολόκληρη τη διαδικασία επιλογής: αναπαράσταση, αρχιτεκτονική, εξαγωγείς, προδιαγραφές και βαθμονόμηση επιλέγονται στις εσωτερικές ομάδες. Η pooled επίδοση εκτός εκπαίδευσης αναφέρεται μαζί με τις ανά γύρο επιδόσεις, επειδή διαφορετική βαθμονόμηση επηρεάζει την κοινή κατάταξη. Η επιβεβαίωση της Α2 χρησιμοποιεί χωριστό κλειστό σύνολο ανεξάρτητων ομάδων, εφόσον η απογραφή επιτρέπει επαρκή κάλυψη· διαφορετικά η μελέτη παραμένει εκτιμητική. Το σύνολο αυτό εξετάζεται μία φορά, αφού οριστικοποιηθούν όλες οι επιλογές. Κάθε τεκμήριο έχει ρόλο ορισμού, ανάπτυξης ή ανεξάρτητης επικύρωσης, και εμφάνιση που διαμόρφωσε προδιαγραφή δεν την επιβεβαιώνει. Πέντε κοινές αρχικοποιήσεις ελέγχουν την εκπαίδευση, ενώ η αβεβαιότητα εκτιμάται σε επίπεδο ανεξάρτητων ομάδων. Τα ουσιώδη μεγέθη επίδρασης, η ισχύς και οι προϋποθέσεις έγκυρης επαναδειγματοληψίας καθορίζονται από τον πραγματικό πιλότο πριν από τις κύριες συγκρίσεις (Kapoor και Narayanan, 2023).

% ============================================================= 5
\section{Μ1: κοινή οριοθέτηση και αναγνώριση}\label{sec:m1}
Η Μ1 απαντά στο πρώτο ερώτημα: πώς πρέπει να διαβαστεί η γραφή ώστε όρια και οικογένεια να εκτιμώνται μαζί. Στον σχεδιασμό του συστήματος η μονάδα εισόδου και η μονάδα-στόχος διαχωρίζονται. Η είσοδος είναι ο τρόπος με τον οποίο διαβάζεται η γραφή, ενώ η θέση είναι αυτό που αναζητείται σε κάθε ανάγνωση: ένα κύλισμα μπορεί να διαβαστεί ως σειρά γραφικών σημείων, ατομικών μελωδικών κινήσεων ή σύνθετων νευματικών ομάδων, και σε κάθε περίπτωση το μοντέλο οφείλει να επιστρέψει το ίδιο συνεχές τμήμα με την ετικέτα της οικογένειάς του.

\paragraph{Δώδεκα αναγνώσεις της γραφής}
Συγκρίνονται τρεις μονάδες εισόδου $U=\{u_1,u_2,u_3\}$. Το γραφικό σημείο $u_1$ διατηρεί το σημάδι αλλά κάποτε συνενώνει περισσότερες κινήσεις· το ατομικό γεγονός $u_2$ δίνει κοινή στοιχειώδη κλίμακα αλλά χάνει την ομαδοποίηση του γραφέα· η σύνθετη νευματική μονάδα $u_3=f_{\mathrm{group}}(N)$ διατηρεί την ομαδοποίηση με ντετερμινιστικό κανόνα $f_{\mathrm{group}}$, που υπολογίζεται μόνο από το σημειογραφικό στρώμα $N$, ώστε να μη μεταφέρει πληροφορία από άλλα στρώματα ή από ετικέτες. Κάθε μονάδα συνδυάζεται με τέσσερα σωρευτικά σύνολα γνωρισμάτων $\Phi=\{N,\,N{+}T,\,N{+}T{+}M,\,N{+}T{+}M{+}L\}$: σημειογραφικά ($N$), χρονικά ($T$), τροπικά ($M$: ήχος, μαρτυρία, φθορά) και κειμενικά-προσωδιακά ($L$: συλλαβή, λέξη, τόνος, μελισματική πυκνότητα). Οι δώδεκα συνθήκες του πλέγματος $U\times\Phi$ συγκρίνονται στις ίδιες τρεις αρχιτεκτονικές με ίσο προϋπολογισμό ρύθμισης, ώστε να εκτιμάται η συμβολή κάθε στρώματος.

\paragraph{Κοινή και διαδοχική πρόβλεψη}
Η Μ1 συγκρίνει τρεις ήδη σχεδιασμένες λύσεις: \textbf{A}, joint semi-CRF που βαθμολογεί κάθε υποψήφιο τμήμα $[s,e]$ μαζί με την οικογένεια· \textbf{B}, boundary-CRF και κατόπιν οικογενειακό ταξινομητή· \textbf{C}, semi-CRF κατάτμησης χωρίς οικογενειακή ανάδραση και τον ίδιο ταξινομητή. Έχουν κοινό τύπο κωδικοποιητή, διαχωρισμούς, διαθέσιμη εποπτεία και προϋπολογισμό. Στις B/C η οικογενειακή απώλεια δεν εκπαιδεύει τον κατατμητή. Η κοινή αξιολόγηση μετρά έτσι αν η ταυτότητα βοηθά να βρεθούν τα όρια. Η A εκπαιδεύεται από μερικές επισημειώσεις και αθροίζει την πιθανοφάνεια των αποδεκτών ζευγών ορίων $\mathcal B_x$. Οι βαθμολογίες χρησιμοποιούν συμβολικά γνωρίσματα ή μικρό νευρωνικό κωδικοποιητή (Kong κ.ά., 2016), με μικρό Transformer--CRF/Transformer--semi-CRF ως σύγχρονη αναφορά. Η οικογένεια εκτός γνωστού λεξικού ($\varnothing_F$) διακρίνεται από την αποχή λόγω ανεπαρκών τεκμηρίων.

Ο event-level αυτοπαλίνδρομος Transformer $T_{\mathrm{evt}}$ εκπαιδεύεται στο ίδιο σώμα και προβλέπει ανά γεγονός ετικέτες αρχής, εσωτερικού ή εκτός θέσης, μαζί με την οικογένεια, με κοινό τελικό $F1_{\mathrm{seg}}$. Ελέγχει έτσι τη χρησιμότητα ρητής τμηματικής βαθμολόγησης. Η εκδοχή του με semi-CRF κεφαλή εξετάζει χωριστά την αναπαράσταση, αφού επαναφέρει το τμήμα. Ο Pitman--Yor προσφέρει μη εποπτευόμενη στατιστική κατάτμηση (Lanz και Hajič, 2025)· για οικογενειακό $F1_{\mathrm{seg}}$ συνδέεται με τον ίδιο εποπτευμένο ταξινομητή, με αυτή την πρόσθετη εποπτεία ρητή.

\paragraph{Γενικές ενδείξεις ορίων}
Ένα κρίσιμο ερώτημα είναι πόσα από τα όρια των θέσεων εξηγούν ήδη οι ενδείξεις που χρησιμοποιεί η μελωδική κατάτμηση χωρίς καμία γνώση των οικογενειών. Τρεις οικογένειες τέτοιων ενδείξεων υπολογίζονται από τη συμβολική είσοδο, με ρυθμίσεις που ορίζονται στα δεδομένα ανάπτυξης: τοπικές ασυνέχειες κίνησης και διάρκειας (LBDM· Cambouropoulos, 2001a), πληροφοριακή απροσδοκησία του επόμενου γεγονότος (IDyOM· Pearce κ.ά., 2010) και επανάληψη μελωδικών δομών (Cambouropoulos, 2006· Conklin και Anagnostopoulou, 2006). Οι ενδείξεις χρησιμοποιούνται με δύο τρόπους. Ως αυτόνομοι κατατμητές, δηλαδή ως μοντέλα που προτείνουν μόνα τους όρια, μαζί με έναν μαθημένο κατατμητή φράσεων με συνελικτικό δίκτυο και επίπεδο CRF (Guan κ.ά., 2018), δίνουν τα μοντέλα αναφοράς για τα όρια. Ως πρόσθετα γνωρίσματα του semi-CRF, με παράλειψη κάθε οικογένειας ενδείξεων χωριστά, δείχνουν αν η γενική μελωδική άρθρωση και η μουσικολογική ταυτότητα είναι συμπληρωματικές.

Η ανάγκη της κοινής εκτίμησης φαίνεται στα παραδείγματα της §\ref{sec:unit}. Όταν μια εναρκτήρια θέση καλύπτει τις τρεις συλλαβές «Οἱ τὰ Χε-», μια τοπική ασυνέχεια ή μια αλλαγή συλλαβής μπορεί να προτείνει όριο μέσα στη θέση· όταν αντίθετα μία συλλαβή αναπτύσσεται σε περισσότερες διαδοχικές θέσεις (Καραγκούνης, 2003, σσ.~207, 209), το όριο πρέπει να τοποθετηθεί μέσα στο μέλισμα, χωρίς καμία κειμενική ένδειξη. Και στις δύο περιπτώσεις το σωστό όριο εξαρτάται από το ποια θέση αναγνωρίζεται. Στον πιλότο τέτοια χωρία καταγράφονται με πλήρεις σημειογραφικές παραπομπές και γίνονται μέρος του συνόλου δύσκολων περιπτώσεων.

\paragraph{Μετρικές και επιλογή αναπαράστασης}
Το $F1_{\mathrm{seg}}$ συνδυάζει ακρίβεια και ανάκληση τμημάτων: κάθε πρόβλεψη αντιστοιχίζεται το πολύ σε μία πραγματική εμφάνιση, με σωστή οικογένεια και ακριβή όρια σε κάποιο ζεύγος του $\mathcal B_x$. Αναφέρονται επίσης $F1_{\mathrm{bnd}}$ ορίων με ανοχή ενός ατομικού γεγονότος, μέση ανάκληση οικογενειών $\mathrm{Rec}_{\mathrm{fam}}$, και $\mathrm{AP}_{\mathrm{fam}}$: one-versus-rest κατάταξη εμφανίσεων για κάθε οικογένεια με γνωστά όρια και κατόπιν macro μέσος. Η κάθε εμφάνιση με εναλλακτά όρια μετρά μία φορά, με μέσο score στα αποδεκτά όριά της. Ισοβαθμίες εισάγονται μαζί στην AP. Κάλυψη χωρίς αποχή, κίνδυνος ανά κάλυψη και αναγνώριση άγνωστων οικογενειών αξιολογούνται χωριστά, με κατώφλια από την ανάπτυξη.

Πρώτα επιλέγεται αναπαράσταση με ισόβαρη αποτίμηση A/B/C, ώστε η A να μην καθορίζει μόνη της το πεδίο σύγκρισης:
\[
\Rs=\arg\max_{R:\,\min_{a\in\{A,B,C\}}\widehat{\mathrm{Rec}}_{\mathrm{fam}}(R,a)\ge r_{\min}}
\frac13\sum_{a\in\{A,B,C\}}\widehat{F1}_{\mathrm{seg}}(R,a).
\]
Έπειτα επιλέγεται αρχιτεκτονική στην ίδια $\Rs$, με ίδιο κριτήριο και κατώφλι. Και τα δύο στάδια ανήκουν στην εσωτερική ανάπτυξη και αξιολογούνται σε εξωτερικές ομάδες. Το $r_{\min}$ οριστικοποιείται στον πιλότο· αν καμία αναπαράσταση το περνά, αναφέρεται η ανεπάρκεια. Μία δεύτερη αναπαράσταση και τα ανά αρχιτεκτονική βέλτιστα ελέγχουν την ευαισθησία. Η τελική $\Rs$ οριστικοποιείται στον μήνα 22 για την Α2 και τη Μ2.

% ============================================================= 6
\section{Α2: περιεχόμενο και μορφή της μουσικολογικής γνώσης}\label{sec:a2}
Η Α2 μετατρέπει τη μουσικολογική προδιαγραφή σε ελεγχόμενη πειραματική μεταβλητή: μετρά την επίδραση του έγκυρου περιεχομένου και το πρόσθετο όφελος της λογικής σύνθεσής του. Η σχέση με τη μελωδική αφαίρεση και τον σχηματισμό κατηγοριών είναι άμεση: η προδιαγραφή δηλώνει ποια γνωρίσματα και ποιες σχέσεις τους είναι διαγνωστικά για κάθε οικογένεια (Cambouropoulos, 2001b· Volk και van~Kranenburg, 2012).

\keybox{\textbf{Σχηματικό παράδειγμα.} Ο κανόνας $z_1\wedge(z_2\vee z_3)\wedge\neg z_4$ μπορεί να δηλώνει απαιτούμενη κατάληξη, δύο εναλλακτικές εσωτερικές αναπτύξεις και αποκλειόμενη τροπική μεταβολή. Η συγκεκριμένη Boolean απόφαση αναπαράγεται και από το γραμμικό κατώφλι $3z_1+z_2+z_3-6z_4\ge4$. Αντίθετα, δύο ανεξάρτητες εναλλακτικές ομάδες, $(a\vee b)\wedge(c\vee d)$, δεν είναι γενικά γραμμικά διαχωρίσιμες. Τα παραδείγματα καθορίζουν ελέγξιμα όρια της σύνθεσης· η αξία των πραγματικών προδιαγραφών μετριέται στο corpus.}

\paragraph{Κοινή πληροφορία και ενεργή παρέμβαση}
Οι εξαγωγείς υπολογίζουν τις έννοιες από τη συμβολική είσοδο για κάθε υποψήφιο τμήμα και κάθε υποψήφια οικογένεια, πριν χρησιμοποιηθεί η ετικέτα. Οι ειδικοί ελέγχουν αξιοπιστία και ελλείψεις στον πιλότο. Η κύρια σύγκριση χρησιμοποιεί παρατηρημένες Boolean τιμές και χωριστή ένδειξη άγνωστης τιμής. Και οι δύο τρόποι χρήσης λαμβάνουν πλήρη πίνακα λεκτικών, ρόλων και εναλλακτικών ομάδων. Η επίπεδη χρήση $F$ συνθέτει προσθετικά τις τρεις προκαθορισμένες συνόψεις ρόλων $Q_f^K$, ενώ η δομημένη $D$ προσθέτει δείκτη παραβίασης της συνολικής λογικής προδιαγραφής. Έτσι η τελική λογική απόφαση αποτελεί διαφορά τελεστή πάνω στην ίδια είσοδο.

Στο κύριο πείραμα η κοινή βασική βαθμολογία $b_{\hat\psi}(s,f)$ μαθαίνεται χωρίς προδιαγραφές στη Μ1 και διατηρείται σταθερή. Με $Q_f^K$ τους τρεις προδηλωμένους μέσους ρόλων και $v_f^K=1[\phi_f^K=0]$:
\[
\ell_{F,K}=b_{\hat\psi}+w^\top Q_f^K,\qquad
\ell_{D,K}=b_{\hat\psi}+w^\top Q_f^K-\lambda v_f^K,
\quad w_r\ge0.1,\ \sum_rw_r=1,\ \lambda\in[0,2].
\]
Τα $w$ είναι κοινά σε όλες τις οικογένειες και εκπαιδεύονται χωριστά ανά συνθήκη, με ίδια διαδικασία. Η οριστικοποιημένη βάση και η σταθερή κλίμακα αποτρέπουν ακύρωση της παρέμβασης από ελεύθερο νευρωνικό κλάδο. Η ποινή εισέρχεται στην πιθανοτική κανονικοποίηση του semi-CRF κατά εκπαίδευση και πρόβλεψη. Τριτιμή λογική διακρίνει τεκμηριωμένη παραβίαση από άγνωστη αλήθεια· η δεύτερη δεν τιμωρείται στην κύρια ανάλυση. Η άγνωστη οικογένεια δεν έχει κανόνα, αλλά συμμετέχει στον ίδιο παρονομαστή. Η πλήρης νευρωνική επανεκπαίδευση εξετάζεται χωριστά ως πρακτική ευαισθησία.

\paragraph{Έξι συνθήκες και δύο κύρια μεγέθη}
Οι τέσσερις συνθήκες $F/D\times V/P$ συνδυάζουν έγκυρες προδιαγραφές με δέκα ελεγχόμενες διαταραχές. Ο γεννήτορας διατηρεί λογικό σκελετό, επικαλύψεις οικογενειών, τροπική εφαρμοσιμότητα και αντιστοιχισμένες συχνότητες/ελλείψεις. Κάθε διαταραχή απαιτεί μη συνωνυμία, ικανοποιησιμότητα, πραγματική ενεργοποίηση και αριθμητικό πιστοποιητικό μη ισοδυναμίας των επιτρεπόμενων κλάσεων F και D. Οι διαταραχές επιλέγονται χωρίς φίλτρο επίδοσης στις ετικέτες και διασταυρώνονται με πέντε κοινά seeds. Συμπληρώνουν τον σχεδιασμό ο $B_K$, με τις ίδιες έννοιες και ελεύθερο γραμμικό κεφάλι χωρίς προδιαγραφή, και ο $B^\star$, με αντιπαραθετική μάθηση SupCon, πρωτότυπα οικογενειών και αντιστοιχισμένη εποπτεία εννοιών. Η μεγαλύτερη χωρητικότητά του δηλώνεται ως εξωτερική ανταγωνιστική σύγκριση.

Αν $S(A,K)$ είναι η μέση $\mathrm{AP}_{\mathrm{fam}}$ στα πέντε seeds και $\bar S(A,P)$ ο μέσος των δέκα διαταραχών, τα κύρια μεγέθη είναι
\[
\begin{aligned}
\Delta_{\mathrm{know}}&=\tfrac12\{S(D,V)-\bar S(D,P)+S(F,V)-\bar S(F,P)\},\\
\Delta_{\mathrm{logic},V}&=S(D,V)-S(F,V).
\end{aligned}
\]
Το πρώτο εκτιμά την επίδραση του έγκυρου περιεχομένου κατά μέσο όρο στους δύο τρόπους χρήσης· το δεύτερο το πρόσθετο όφελος της συγκεκριμένης δομημένης λογικής σύνθεσης έναντι της προσθετικής σύνοψης της ίδιας έγκυρης γνώσης. Η μέση δομική επίδραση και η αλληλεπίδραση παραμένουν δευτερεύουσες, μαζί με $F1_{\mathrm{seg}}$, δύσκολες περιπτώσεις και αποτελέσματα ανά οικογένεια. Θετική επίδραση έναντι P ερμηνεύεται μαζί με τη σύγκριση έναντι $B_K$.

Οι δύο κύριοι μονόπλευροι έλεγχοι χρησιμοποιούν Holm με οικογενειακό $\alpha=0.025$. Αν απορριφθούν και οι δύο μηδενικές, ανοίγει δεύτερη οικογένεια Holm για υπεροχή της DV έναντι $B_K$ και $B^\star$. Η επιβεβαίωση απαιτεί ανεξάρτητες ομάδες και ελεγμένη συμπεριφορά επαναδειγματοληψίας της ίδιας pooled AP· sign-flip της μέσης log-πιθανότητας ελέγχει διαφορετικό μέγεθος. Με ανεπαρκή δεδομένα παραδίδεται εκτιμητική Α2 και συνεχίζεται ο υπόλοιπος κορμός. Το συνοδευτικό τεχνικό πρωτόκολλο καθορίζει εξισώσεις, generator, στατιστικά gates και πλήρη καταγραφή. Οι μαθηματικοί έλεγχοι πιστοποιούν την κατασκευή· οι επιδόσεις της θα προκύψουν από πραγματικά έργα.

% ============================================================= 7
\section{Μ2: ταυτότητα, πραγμάτωση και ιστορική εποπτεία}\label{sec:repr}
Η Μ1 προσδιορίζει την αναπαράσταση και η Α2 εκτιμά την αξία των διαγνωστικών σχέσεων. Το τρίτο ερώτημα εξετάζει τι μαθαίνει ένας κωδικοποιητής όταν παρατηρεί πολλές πραγματώσεις της ίδιας οικογένειας, και αν κερδίζει σε αναγνώριση χωρίς να σβήνει τη μουσικά σημαντική παραλλαγή.

\paragraph{Ο τμηματικός κωδικοποιητής}
Ο κωδικοποιητής δέχεται την $\Rs$ και υπολογίζει για κάθε γεγονός μια τοπική κατάσταση $h_t$, δηλαδή εσωτερική περιγραφή του γεγονότος μαζί με τα συμφραζόμενά του. Από τις καταστάσεις ενός τμήματος παράγονται δύο αναπαραστάσεις: η $z_{\mathrm{id}}$, για την ταυτότητα της οικογένειας, και η $z_{\mathrm{real}}$, για την πληροφορία της συγκεκριμένης πραγμάτωσης, που διατηρείται με βοηθητικό στόχο ανακατασκευής. Η διάκριση των δύο ρόλων ελέγχεται εμπειρικά. Η σχεσιακή απώλεια εφαρμόζεται στο $z_{\mathrm{id}}$, αλλά η κλίση της φθάνει στον κοινό κωδικοποιητή. Σε αχρησιμοποίητες ιστορικές ακμές ελέγχονται χωριστά η έλξη στο $z_{\mathrm{id}}$, η μεταβολή της τοπικής γεωμετρίας $h$ και η διατήρηση πραγματωτικών γνωρισμάτων με το $\Qreal$. Ο κοινός στόχος ανακατασκευής δεν εγγυάται αμετάβλητο $h$. Η αυτοεπιβλεπόμενη προεκπαίδευση, μια αρχική εκπαίδευση χωρίς ετικέτες στην οποία ο κωδικοποιητής μαθαίνει να συμπληρώνει κρυμμένα συμβολικά γεγονότα (Zeng κ.ά., 2021), αφορά μόνο τον κωδικοποιητή: μία πολιτική απόκρυψης, γεγονότων, τμημάτων ή ολόκληρων συλλαβών, επιλέγεται στα δεδομένα ανάπτυξης και εφαρμόζεται ισότιμα στις B1 και B2, και η συμβολή της ελέγχεται συγκρίνοντας τον κωδικοποιητή με και χωρίς προεκπαίδευση.

\paragraph{Τρεις κλιμακωτές συνθήκες}
Η \textbf{B0} είναι το semi-CRF της Μ1 με την $\Rs$ και συμβολικά γνωρίσματα, χωρίς μάθηση αναπαράστασης, και δείχνει πόσα εξηγεί ήδη η επιφάνεια της γραφής. Η \textbf{B1} προσθέτει τον προεκπαιδευμένο κωδικοποιητή με εποπτεία οικογένειας και ανακατασκευή, χωρίς σχεσιακό όρο. Η \textbf{B2} διατηρεί την ίδια αρχιτεκτονική, αρχικοποίηση και εκπαίδευση και προσθέτει έναν όρο που έλκει στο $z_{\mathrm{id}}$ ζεύγη πραγματικών εμφανίσεων της ίδιας οικογένειας από διαφορετικά έργα:
\[
\mathcal L_{\mathrm{B2}}=\mathcal L_{\mathrm{seg}}+\lambda_f\,\mathcal L_{\mathrm{fam}}+\lambda_r\,\mathcal L_{\mathrm{rec}}+\lambda_p\,\mathcal L_{\mathrm{pair}},\qquad \lambda_p=0\ \text{στη B1}.
\]
Στην έκφραση αυτή, $\mathcal L_{\mathrm{seg}}$ είναι η απώλεια του semi-CRF για όρια και οικογένεια, $\mathcal L_{\mathrm{fam}}$ η απώλεια της ταξινόμησης στην οικογένεια, $\mathcal L_{\mathrm{rec}}$ η απώλεια της ανακατασκευής, $\mathcal L_{\mathrm{pair}}$ ο όρος έλξης θετικών ζευγών, που διακρίνεται από την πλήρη SupCon με αρνητικά του $B^\star$ (Khosla κ.ά., 2020), και τα $\lambda$ τα βάρη τους, που επιλέγονται στα δεδομένα ανάπτυξης. Κύριο μέγεθος είναι το $\Delta_{\mathrm{famrel}}=\mathrm{AP}_{\mathrm{fam}}(\mathrm{B2})-\mathrm{AP}_{\mathrm{fam}}(\mathrm{B1})$ σε όλες τις εμφανίσεις εκτός εκπαίδευσης, με τις δύσκολες περιπτώσεις ως δευτερεύον υποσύνολο. Επειδή ένας κωδικοποιητής μπορεί να πετύχει καλύτερη αναγνώριση συμπιέζοντας κάθε διαφορά, το $\Delta_{\mathrm{famrel}}$ ερμηνεύεται πάντα μαζί με το $\Qreal$.

\paragraph{Πληροφορία πραγμάτωσης και συντομεύσεις}
Το $\Qreal$ μετρά πόση πληροφορία της συγκεκριμένης πραγμάτωσης μένει προσβάσιμη στις αναπαραστάσεις των εμφανίσεων εκτός εκπαίδευσης, με τρεις προκαθορισμένες δοκιμές: ανακατασκευή αποκρυμμένων συμβολικών γνωρισμάτων, διάκριση δύο πραγματώσεων της ίδιας οικογένειας ως προς τη μελωδική και χρονική τους διάρθρωση, και πρόβλεψη τοπικών ιδιοτήτων, όπως η έκταση και η συλλαβική κατανομή, από απλό γραμμικό ανιχνευτή, δηλαδή ένα μοντέλο περιορισμένης χωρητικότητας που ελέγχει τι μπορεί να διαβαστεί από την αναπαράσταση. Οι δοκιμές εφαρμόζονται με τον ίδιο τρόπο σε όλες τις συνθήκες, και για τη B2 δηλώνεται εκ των προτέρων ένα όριο μη κατωτερότητας, δηλαδή η μέγιστη απώλεια πληροφορίας έναντι της B1 που θεωρείται αποδεκτή. Ισχυρή ομοιογένεια του $z_{\mathrm{id}}$ με απώλεια της πραγμάτωσης αναφέρεται ως χωριστό εύρημα, ακόμη κι αν η αναγνώριση βελτιώνεται. Επειδή εμφανίσεις της ίδιας οικογένειας συχνά μοιράζονται κείμενο, καταληκτική βαθμίδα ή μήκος, η σύγκριση επαναλαμβάνεται χωρίς το κειμενικό στρώμα $L$ στον κλάδο ταυτότητας, και οι έλεγχοι οικογενειακής διάκρισης περιλαμβάνουν δύσκολα αρνητικά παραδείγματα υψηλής επιφανειακής ομοιότητας, χωρίς προσθήκη αρνητικών στον όρο έλξης $\mathcal L_{\mathrm{pair}}$ (Sturm, 2014).

\subsection{Ιστορικό σκέλος της Μ2: B2--EISE}\label{sec:historicpriority}
Η τεκμηρίωση συν-εξηγήσεων ανήκει στο προπαρασκευαστικό έργο της Μ2. Πιλοτική εκπαίδευση ενεργοποιείται μόνο με επιλέξιμη τράπεζα. Η \textbf{EISE} (Exegesis-Invariant Segmental Encoder) χρησιμοποιεί την αρχιτεκτονική και την απώλεια της B2, αντικαθιστώντας τα ζεύγη διαφορετικών έργων κοινής οικογένειας με ανεξάρτητα επαληθευμένες συν-εξηγήσεις του ίδιου χωρίου. Αντιστοιχίζονται έκθεση, βαθμοί, οικογένεια, προέλευση, επιφανειακή απόσταση και προϋπολογισμός· η πραγματική εφικτότητα ελέγχεται πριν από την εκπαίδευση. Η ίδια ιστορική ταυτότητα του χωρίου παραμένει το διαφορετικό χαρακτηριστικό της εποπτείας.

Οι ίδιες ανεξάρτητες ιστορικές εμφανίσεις ευθυγραμμίζονται σε κοινές μουσικές περιοχές $u$, χωρίς χρήση εκτελεστικού χρονισμού. Με μέση τιμή $\mu_{m,\mathrm{tr}}$ και συνδιακύμανση $\Sigma_{m,\mathrm{tr}}$ μόνο από την εκπαίδευση ορίζονται
\[
\begin{gathered}
a_m=\operatorname{tr}(\Sigma_{m,\mathrm{tr}})/d_h>0,\qquad
W_m=(\Sigma_{m,\mathrm{tr}}+0.01a_m I)^{-1/2},\\
\widetilde h_m(u)=W_m(h_m(u)-\mu_{m,\mathrm{tr}}),\qquad
\widehat\tau_{f,m,\mathrm{hist}}(u)=
\sqrt{\operatorname{tr}\widehat\Sigma_{f,m}(u)/d_h}.
\end{gathered}
\]
Η $\widehat\Sigma$ χρησιμοποιεί κοινά βάρη, ίσα ανά ανεξάρτητη ομάδα πηγών και κατόπιν ανά χωρίο και μάρτυρα. Η Μ4-min παρέχει το ανεξάρτητο $\Gamma_f(u)$, από ηχογραφήσεις όπως το σώμα ΔΑΜΑΣΚΗΝΟΣ (Delviniotis και Kouroupetroglou, 2011): διασπορά κανονικοποιημένου λογαριθμικού χρονισμού μεταξύ ψαλτών στα ίδια χωρία. Η πραγματική τομή ηχογραφήσεων και ιστορικών πηγών ελέγχεται στην απογραφή. Ορίζονται
\[
S_m=\frac{1}{|\mathcal F_{\mathrm{eval}}|JR}\sum_{f\in\mathcal F_{\mathrm{eval}}}\sum_{j=1}^{J}\sum_{r=1}^{R}\operatorname{Spearman}_{u}\bigl(\widehat\tau_{f,mjr,\mathrm{hist}}(u),\Gamma_f(u)\bigr),\quad \dhist=S_{\mathrm{EISE}}-S_{\mathrm{B2}},
\]
όπου $j$ δηλώνει προκαθορισμένο γράφημα ζευγοποίησης B2 και, στην EISE, την αντίστοιχη εκπαιδευτική επανάληψη· το $r$ δηλώνει κοινή αρχικοποίηση. Η κύρια σύγκριση αφορά την οριστικοποιημένη τράπεζα $J$. Θετικό $\dhist$ σημαίνει υψηλότερη συμφωνία κατάταξης της EISE· αναφέρονται και τα δύο $S_m$. Η συμπίεση δεν επιβάλλει πρόσημο, ούτε το $\dhist$ απομονώνει αιτιωδώς την ιστορική προέλευση από την κοινή ταυτότητα του χωρίου. Το τοπικό ιστορικό προφίλ περιλαμβάνει διαφοροποίηση μεταξύ χωρίων και μεταξύ μαρτύρων, που εξετάζονται χωριστά διαγνωστικά. Σταθερά προφίλ ή ανύπαρκτο κοινό πεδίο συνεπάγονται αποχή, όχι μηδενική συσχέτιση.

Η μία κύρια ιστορική αντίθεση ελέγχεται αμφίπλευρα, με κοινή επαναδειγματοληψία πλήρων προφίλ, ομάδων πηγών και ψαλτών και διασταυρούμενες επαναλήψεις· η διαδικασία πιστοποιείται πρώτα με προσομοίωση της πραγματικής δομής εξαρτήσεων. Επιβεβαίωση απαιτεί $G_{\mathrm{pair}}$, $G_{\mathrm{perf}}$ και ελεγμένη κάλυψη/ισχύ της συγκεκριμένης διαδικασίας. Ο μήνας 6 εγκρίνει την εφικτότητα· η τελική βαθμονόμηση ακολουθεί την επιλογή $\Rs$, πριν από το ιστορικό σύνολο ελέγχου. Η άμεση ιστορική--εκτελεστική συσχέτιση $\rho_A$, όπου ορίζεται, παραμένει χωριστό περιγραφικό αποτέλεσμα. Με ανεπαρκή ακρίβεια αναφέρεται μόνο εκτιμητικό αποτέλεσμα, ενώ με μη ορισμένο προφίλ ή χωρίς εκτελεστικό πεδίο το $\dhist$ δεν εκτιμάται. Οι ακριβείς κανόνες δίνονται στο συνοδευτικό πρωτόκολλο Μ2.

% ============================================================= 8
\section{Ανεξάρτητη μουσικολογική αξιολόγηση}\label{sec:validation}
Όλες οι παραπάνω συγκρίσεις κρίνονται τελικά απέναντι σε ανεξάρτητη μουσικολογική κρίση. Στη Μ3a, δύο ειδικοί με διαφορετική διδακτική ή ερευνητική προέλευση, που δεν συνέταξαν τις προδιαγραφές, κρίνουν πρώτα τυφλά τα όρια και την ένταξη νέων χωρίων (μέλος, μη μέλος, ανεπίλυτο), χωρίς πρόσβαση στις προδιαγραφές ή στις προβλέψεις. Αυτές οι ανεξάρτητες κρίσεις δίνουν το κριτήριο αξιολόγησης. Σε δεύτερη φάση εξετάζουν την εφαρμοσιμότητα των προδιαγραφών και τις διαγνωστικές σχέσεις τους. Οι κρίσεις ακολουθούν σταθερό εγχειρίδιο, και οι διαφωνίες ταξινομούνται κατά αιτία: όριο, επίπεδο ανάλυσης, τροπικό ζήτημα, γνώρισμα, πηγή ή ιστορική ονομασία. Οι διαφωνίες αποτελούν οι ίδιες εύρημα, αφού δείχνουν πού τα όρια είναι πραγματικά αμφίσημα και ποιες μορφολογικές διακρίσεις αμφισβητούνται, και η επίλυσή τους διατηρεί καταγεγραμμένη την αρχική διαφωνία. Η πλήρης τυφλή επικύρωση γίνεται στους μήνες 12--20, αφού οι προδιαγραφές πάρουν την τελική τους έκδοση, και τροφοδοτεί το σύνολο επιβεβαίωσης.

\paragraph{Ρόλος του ερευνητή και δεοντολογία}
Οι προδιαγραφές διατυπώνονται από τον ερευνητή, και γι' αυτό οι δικές του κρίσεις καταχωρίζονται μόνο ως υλικό ορισμού και ανάπτυξης. Η τελική εφαρμοσιμότητα κρίνεται από τους ανεξάρτητους ειδικούς, και κάθε αλλαγή προδιαγραφής καταγράφεται με αριθμό έκδοσης. Τα πρωτόκολλα με ανθρώπινους συμμετέχοντες υποβάλλονται στην Επιτροπή Ηθικής και Δεοντολογίας της Έρευνας του Ελληνικού Μεσογειακού Πανεπιστημίου (ΕΛΜΕΠΑ) ήδη στο ΠΕ0, με ενημερωμένη συναίνεση και ελάχιστη συλλογή προσωπικών δεδομένων.

\paragraph{Υπολογιστικοί πόροι}
Τα μοντέλα είναι μικρά και εκπαιδεύονται σε έναν επιταχυντή γραφικών μεσαίας κατηγορίας ή σε αντίστοιχη πανεπιστημιακή υποδομή. Ο πιλότος μετρά χρόνο και μνήμη. Η Α2 προβλέπει 120 εκπαιδεύσεις ανά τελικό διαχωρισμό (κυρίως μικρών σταθερών κεφαλών). Η πλήρης Μ2b με $J=10$ γραφήματα και $R=5$ αρχικοποιήσεις προβλέπει έως $2JR=100$ εκπαιδεύσεις B2/EISE, πέρα από B1, προεκπαίδευση και ρύθμιση· ελέγχεται πιλοτικά αν αρκεί οικονομικότερη ισοδύναμη διασταύρωση.

% ============================================================= 9
\section{Χρονοδιάγραμμα, δημοσιεύσεις και κίνδυνοι}\label{sec:plan}
Το πρόγραμμα οργανώνεται γύρω από τέσσερα σημεία δέσμευσης: τα αρχικά κριτήρια επάρκειας στον μήνα 6, η τελική έκδοση των προδιαγραφών και ο ορισμός των συνόλων επιβεβαίωσης στον μήνα 12, η επιλογή της $\Rs$ στον μήνα 22 και ο τελικός έλεγχος μόνο αφού οριστικοποιηθούν οι συγκρίσεις (Πίνακας~\ref{tab:timeline}). Η εξειδικευμένη στατιστική βαθμονόμηση της Μ2b ακολουθεί την επιλογή $\Rs$, πριν από το ιστορικό σύνολο ελέγχου.

\begin{table}[htbp]
\caption{Χρονοδιάγραμμα 48 μηνών, σημεία δέσμευσης και δημοσιεύσεις}\label{tab:timeline}
\footnotesize
\begin{tabularx}{\linewidth}{@{}L{1.3cm}YL{5.2cm}@{}}
\toprule
\textbf{Μήνες} & \textbf{Εργασία} & \textbf{Παραδοτέο ή σημείο απόφασης}\\\midrule
1--6 & ΠΕ0· πιλότος σε τρεις οικογένειες· πρώτη δοκιμή εξαγωγής εννοιών· απογραφή συν-εξηγήσεων και ηχογραφήσεων & Έκθεση πιλότου· απόφαση για τα τέσσερα βασικά κριτήρια· προκαταρκτικός έλεγχος εφικτότητας $G_{\mathrm{pair}}$, $G_{\mathrm{perf}}$\\
7--12 & Επέκταση κωδικοποίησης και επισημείωσης· προδιαγραφές και εξαγωγείς εννοιών & Τελική έκδοση προδιαγραφών· ορισμός του συνόλου επιβεβαίωσης (μήνας 12)\\
12--20 & Μ3a: πλήρης τυφλή επικύρωση· μοντέλα αναφοράς ορίων & Επικυρωμένο λεξικό· \textbf{άρθρο 1}: εργασία, πρωτόκολλο και σώμα (στόχος μήνας 18--20)\\
13--22 & Μ1: πλέγμα $U\times\Phi$ σε A/B/C, κοινή βάση, επιλογή $\Rs$ & Επιλογή της $\Rs$ (μήνας 22)\\
22--30 & Α2: $2\times2$, $B_K$, $B^\star$ και έλεγχοι μη ισοδυναμίας & \textbf{Άρθρο 2}: μουσικολογική γνώση ως πειραματική μεταβλητή (στόχος μήνας 30)\\
24--38 & Μ2: B0--B2 και $\Qreal$· πιλότος EISE μόνο με επιλέξιμα ζεύγη· τελική βαθμονόμηση πριν από το ιστορικό σύνολο ελέγχου & \textbf{Άρθρο 3}: ταυτότητα, πραγμάτωση και ιστορική εποπτεία· επιβεβαιωτικό $\dhist$ μόνο με επαρκές κοινό πεδίο και ακρίβεια\\
39--48 & Σύνθεση, αναπαραγωγιμότητα, συγγραφή & Διατριβή· δημόσιο υλικό εντός αδειών\\\bottomrule
\end{tabularx}
\end{table}

Τα άρθρα στοχεύουν στο ετήσιο συνέδριο της International Society for Music Information Retrieval (ISMIR), στο περιοδικό της \emph{Transactions of the International Society for Music Information Retrieval} και σε συγγενή περιοδικά. Αν οι άδειες δεν επιτρέψουν δημόσια διάθεση των μεταγραφών, το πρώτο άρθρο δημοσιεύει την εργασία, το πρωτόκολλο, τις επισημειώσεις και τα μεταδεδομένα. Οι κύριοι κίνδυνοι αφορούν την κάλυψη και την αξιοπιστία. Η περιορισμένη ψηφιακή διαθεσιμότητα αντιμετωπίζεται με την πρώιμη απογραφή και με περιορισμό του ρεπερτορίου· η ασυμφωνία των ειδικών με αναθεώρηση ή αιτιολογημένη συγχώνευση οικογενειών, με τις αμφίσημες εμφανίσεις να καταγράφονται ως τέτοιες· οι λίγες ανεξάρτητες ομάδες έργων με την ομαδοποιημένη διασταυρούμενη επικύρωση και με περιορισμό των επιβεβαιωτικών ισχυρισμών. Καμία από αυτές τις αποφάσεις δεν αλλάζει εκ των υστέρων σύνολα που έχουν ήδη οριστεί.

% ============================================================= ΠΑΡΑΡΤΗΜΑ
\appendix
\section{Συμπληρωματικοί έλεγχοι και ερευνητικές επεκτάσεις}\label{app:extensions}
Ο ελάχιστος κορμός περιλαμβάνει επισημειωμένα έργα της Νέας Μεθόδου, Μ1, Α2, Μ2a και ανεξάρτητη μουσικολογική επικύρωση Μ3a. Η απογραφή της ιστορικής Μ2b αποτελεί μέρος του προγράμματος· η εκπαίδευση EISE και η επιβεβαιωτική σύγκριση ενεργοποιούνται μόνο με τεκμηριωμένη επάρκεια. Οι παρακάτω συμπληρωματικοί έλεγχοι και επεκτάσεις χρησιμοποιούν την ίδια υποδομή, χωρίς να αποτελούν όλοι προϋπόθεση ολοκλήρωσης.

\subsection{Ιστορική εποπτεία: πρόσθετοι έλεγχοι}
Η §\ref{sec:historicpriority} συμπληρώνεται από πρώιμη απογραφή της επικάλυψης επιφανειακών αποστάσεων, αυστηρότερη αντιστοίχιση κειμένου, $B2_{\mathrm{sim}}$, ανασυνδέσεις βαθμού και αφαίρεση του κειμενικού στρώματος. Αν δεν υπάρχει συγκρίσιμο πεδίο, η χαλαρότερη σύγκριση δηλώνεται διερευνητική και αφορά τη συνολική στρατηγική εποπτείας. Η μετρική ελέγχεται ως προς την παράμετρο σταθεροποίησης και τη σταθερότητα μεταξύ επαναλήψεων. Η απόκρυψη περιοχών στο $z_{\mathrm{id}}$ παραμένει διαγνωστική, χωρίς κανονικοποίηση με το άθροισμα προσημασμένων επιδράσεων. Το πλήρες πρωτόκολλο μεταφέρει τις αρχές των G.3--G.4, E.10--E.11 και F.1--F.5 στη σημερινή απώλεια, με μία κύρια ιστορική αντίθεση. Χωρίς ανεξάρτητο εκτελεστικό προφίλ δεν εκτιμάται $\dhist$. \textit{Προπαρασκευή:} μήνες 1--12· \textit{πιλοτική εκπαίδευση:} μήνες 24--30, εφόσον υπάρχουν επιλέξιμα ζεύγη.

\subsection{Νευροσυμβολική σύνθεση των προδιαγραφών: ECSI}
Η ECSI συνεχίζει το πείραμα Α2. Στην Α2 οι κανόνες των προδιαγραφών εισέρχονται ως γνωρίσματα ή ως όροι ποινής· στην ECSI ένα μοντέλο εκτιμά από κοινού όλες τις διαγνωστικές έννοιες μιας εμφάνισης, και η πιθανότητα ένταξης σε μια οικογένεια υπολογίζεται αθροίζοντας ακριβώς την πιθανότητα όλων των συνδυασμών εννοιών που επιτρέπει η προδιαγραφή στο συγκεκριμένο τροπικό περιβάλλον (Darwiche και Marquis, 2002· Xu κ.ά., 2018). Η πρόσθετη αξία αυτής της ακριβούς σύνθεσης κρίνεται έναντι της απλούστερης εκδοχής που θεωρεί τις έννοιες ανεξάρτητες μεταξύ τους (van~Krieken κ.ά., 2024) και έναντι ενός μοντέλου αναφοράς με αντιπαραθετική μάθηση, στα δύσκολα παραδείγματα και με την προϋπόθεση ότι διατηρείται η πληροφορία της πραγμάτωσης ($\Qreal$). \textit{Προϋπόθεση:} λειτουργική Α2 και επαρκή δύσκολα παραδείγματα, με χωριστό πρωτόκολλο· \textit{νωρίτερη έναρξη:} μήνας 30.

\subsection{Πιθανοτική πραγμάτωση της θέσης: ο τελεστής \texorpdfstring{\Mth}{M-thesis}}
Ο πιθανοτικός τελεστής \Mth\ αντιστρέφει την κατεύθυνση της αναγνώρισης: για δεδομένη οικογένεια θέσεων και μουσικό περιβάλλον εκτιμά την κατανομή των πιθανών μελωδικών πραγματώσεων. Η μελέτη θα μετρήσει πόση πληροφορία προσθέτει η οικογένεια πέρα από το περιβάλλον, σε bits ανά συμβολικό γεγονός, και πώς διακρίνονται οι μορφές που επιτρέπει η προδιαγραφή από εκείνες που προτιμά στην πράξη η παράδοση. Επειδή η ένταξη μιας εμφάνισης στην οικογένεια κρίνεται συνήθως από το ίδιο το μουσικό υλικό που πρέπει να προβλεφθεί, η πρόβλεψη αξιολογείται ως πραγματική πρόβλεψη μόνο όπου η οικογένεια τεκμηριώνεται από πληροφορία έξω από αυτό. \textit{Προϋπόθεση:} αρκετές ανεξάρτητες πραγματώσεις ανά οικογένεια και θετικός έλεγχος του αν η οικογένεια μπορεί να προσδιοριστεί ανεξάρτητα· \textit{νωρίτερη έναρξη:} μήνας 12.

\subsection{Η εξήγηση ως ιστορικός μετασχηματισμός: ATL και ATL-INV}
Η ATL (Attested Transformation Learning) μελετά τεκμηριωμένους μετασχηματισμούς: εξετάζει αν οι κανονικότητες ανάπτυξης και σύμπτυξης ανάμεσα σε επαληθευμένες συν-εξηγήσεις μεταφέρονται από οικογένεια σε οικογένεια. Η ATL-INV (αντίστροφη ιστορική ανάκτηση) εξετάζει αν συγκεκριμένες διακρίσεις της παλαιότερης σημειογραφίας, όπως σημάδια με την ίδια μετροφωνία αλλά διαφορετική μελική ενέργεια, παραμένουν προβλέψιμες από την αναλυτική εξήγηση πέρα από όσα προβλέπουν οι καταγεγραμμένοι κανόνες της εξήγησης. Στόχος είναι η ανάγνωση συγκεκριμένων σωζόμενων μαρτύρων, χωρίς αξίωση ανακατασκευής ενός μοναδικού αρχικού προτύπου. Τα εργαλεία προέρχονται από τη σύγκριση μουσικών ακολουθιών με πράξεις κατάτμησης και συνένωσης (Mongeau και Sankoff, 1990), από την εκμάθηση του κόστους κάθε πράξης από τα ίδια τα δεδομένα (Ristad και Yianilos, 1998) και από τα μοντέλα θορυβώδους καναλιού για την αντίστροφη μεταγραφή (Knight και Graehl, 1998). \textit{Προϋπόθεση:} επαρκείς ανεξάρτητες ιστορικές αντιστοιχίες και παλαιογραφικός έλεγχος των μαρτύρων· \textit{νωρίτερη έναρξη:} μήνας 18.

\subsection{Ακρόαση, εκτέλεση και πλήρη έργα}
Η μουσική σημασία των οικογενειών μπορεί να εξεταστεί και έξω από τη γραφή. Η Μ3b ελέγχει με τριάδες αν η ανθρώπινη κατηγοριοποίηση ακολουθεί τις οικογένειες και η Μ3c με κρίσεις απροσδόκητου αν η πρόβλεψη σε επίπεδο οικογένειας προσθέτει πληροφορία στην πρόβλεψη γεγονότων (Pearce, 2018). Η Μ4-min αποτελεί μόνο το ελάχιστο ανεξάρτητο εκτελεστικό κριτήριο της Μ2b, ενώ η πλήρης Μ4 εξετάζει ευρύτερα τη διαφοροποίηση χρονισμού μεταξύ ψαλτών (Gingras κ.ά., 2016). Η Μ5 εφαρμόζει την αναγνώριση σε πλήρη έργα, εξετάζοντας χωριστά επιλογή οικογενειών, διαδοχή, λειτουργικό περιβάλλον και ενδοοικογενειακή πραγμάτωση, με μελέτη περίπτωσης του Νικολάου Πρωτοψάλτη Σμύρνης (Ιωαννίδης, 2022). \textit{Προϋποθέσεις:} επικυρωμένο λεξικό, πλήρη έργα και συγκρίσιμο σώμα αναφοράς για τη συμβολική Μ5· ηχογραφήσεις και αξιόπιστη στοίχιση για τη Μ4· κατάλληλη δεοντολογική έγκριση για τις ακροαματικές Μ3b--Μ3c· \textit{νωρίτερη έναρξη:} μήνας 24.

\subsection{Συγκριτική εφαρμογή σε άλλη παράδοση}
Εφόσον εξασφαλιστούν συγκρίσιμες συμβολικές πηγές, τεκμηριωμένες μουσικές μονάδες και συνεργασία με ειδικό, η μέθοδος μπορεί να δοκιμαστεί στο ζναμένιο μέλος της ρωσικής ορθόδοξης παράδοσης, όπου επίσης συνοπτικοί σχηματισμοί της γραφής αποδίδονται με αναλυτικότερες μορφές. Η εφαρμογή θα δείξει ποια στοιχεία της αναπαράστασης των οικογενειών και της μάθησης μεταφέρονται και ποια απαιτούν ειδική μουσικολογική αναδιατύπωση. \textit{Προϋπόθεση:} συγκρίσιμο συμβολικό σώμα και ειδικός της παράδοσης· \textit{νωρίτερη έναρξη:} μήνας 36.

\section*{Βιβλιογραφία}
\subsection*{Μουσικολογικές πηγές και μελέτες}
\begin{reflist}
\item Αλεξάνδρου, Μ. (2010). \emph{Εξηγήσεις και μεταγραφές της βυζαντινής μουσικής: Σύντομη εισαγωγή στον προβληματισμό τους}. Θεσσαλονίκη: University Studio Press.
\item Αποστολόπουλος, Θ. Κ. (2002). \emph{Ο Απόστολος Κώνστας ο Χίος και η συμβολή του στη θεωρία της μουσικής τέχνης} (Μελέται 4). Αθήνα: Ίδρυμα Βυζαντινής Μουσικολογίας.
\item Αρβανίτης, Ι. (2010). \emph{Ο ρυθμός των εκκλησιαστικών μελών μέσα από την παλαιογραφική έρευνα και την εξήγηση της παλαιάς σημειογραφίας} (Διδακτορική διατριβή, 2 τόμοι). Ιόνιο Πανεπιστήμιο. \url{https://doi.org/10.12681/eadd/21885}
\item Ιωαννίδης, Χ. (2022). Η σύνθεση του Νικολάου Πρωτοψάλτου Σμύρνης στο Κοινωνικό «Εἰς πᾶσαν τὴν γῆν» σε ήχο πλ. Δ΄ και ο βαθμός εξάρτησής της από την προγενέστερη παράδοση της εκκλησιαστικής μελοποιίας. \emph{Θεολογία}, 93(4), 173--193.
\item Καραγκούνης, Κ. Χ. (2003). \emph{Η Παράδοση και Εξήγηση του Μέλους των Χερουβικών της Βυζαντινής και μετα-Βυζαντινής Μελοποιίας} (Μελέται 7). Αθήνα: Ίδρυμα Βυζαντινής Μουσικολογίας.
\item Καραγκούνης, Κ. Χ. (2005). \emph{Παραλειπόμενα περὶ τοῦ Χερουβικοῦ Ὕμνου}. Βόλος: Γραφικές Τέχνες Αλέξανδρος Ξουράφας.
\item Καράς, Σ. Ι. (1982). \emph{Μέθοδος της Ελληνικής Μουσικής: Θεωρητικόν} (τ. Α΄--Β΄). Αθήνα: Σύλλογος προς Διάδοσιν της Εθνικής Μουσικής.
\item Σπανουδάκης, Δ. Π. (2017). \emph{Σύγχρονα μοντέλα ανάλυσης Βυζαντινής Μουσικής: Το φαινόμενο της μελισματικότητας στα κεκραγάρια} (Διδακτορική διατριβή, 3 τόμοι). Αριστοτέλειο Πανεπιστήμιο Θεσσαλονίκης. \url{https://doi.org/10.12681/eadd/42193}
\item Στάθης, Γ. Θ. (1978). \emph{Η εξήγησις της Παλαιάς Βυζαντινής Σημειογραφίας} (Μελέται 2). Αθήνα: Ίδρυμα Βυζαντινής Μουσικολογίας.
\item Τόμα, Ι. Γ. (2018). \emph{Νικόλαος Γεωργίου, πρωτοψάλτης Σμύρνης: Βίος, ψαλτικό έργο, μουσικολογικές προσεγγίσεις} (Μεταπτυχιακή διπλωματική εργασία). Αριστοτέλειο Πανεπιστήμιο Θεσσαλονίκης.
\item Χρύσανθος ὁ ἐκ Μαδύτων (1832). \emph{Θεωρητικὸν Μέγα τῆς Μουσικῆς}. Τεργέστη.
\item Ψάχος, Κ. Α. (1917/1978). \emph{Η Παρασημαντική της Βυζαντινής Μουσικής} (β΄ επαυξημένη έκδ., επιμ. Γ. Χατζηθεοδώρου). Αθήνα: Διόνυσος.
\item Alexandru, M. (2012). Some thoughts on the so-called great hypostases of the Middle Byzantine notation. \emph{Musicology Papers}, 27(1), 47--58.
\item Alexandru, M., \& Troelsgård, C. (2020). The “Elements of the Papadikê” and modality features in Byzantine chant. \emph{Series Musicologica Balcanica}, 1(1), 169--204.
\item Conomos, D. E. (Ed.). (1985). \emph{The Treatise of Manuel Chrysaphes, the Lampadarios} (Monumenta Musicae Byzantinae, Corpus Scriptorum de Re Musica 2). Vienna: Verlag der Österreichischen Akademie der Wissenschaften.
\item Gheorghiță, N. (2010). Observations on the technique of transcription (\emph{exegesis}) into the New Method of analytical music notation of the Sunday \emph{koinonikon} of the 18th century. \emph{Studii și Cercetări de Istoria Artei. Teatru, Muzică, Cinematografie}, 4(48), 9--26.
\item Gjerdingen, R. O. (2007). \emph{Music in the Galant Style}. Oxford: Oxford University Press.
\item Hannick, C., \& Wolfram, G. (Eds.). (1985). \emph{Gabriel Hieromonachos: Abhandlung über den Kirchengesang} (Monumenta Musicae Byzantinae, Corpus Scriptorum de Re Musica 1). Wien: Verlag der Österreichischen Akademie der Wissenschaften.
\item Rabinovitch, G., \& Carter-Enyi, A. (2024). Melodic organization and sequential ordering of galant schemata: Implications for eighteenth-century European musicianship. \emph{Music Theory Online}, 30(1). \url{https://doi.org/10.30535/mto.30.1.7}
\item St. Anthony’s Greek Orthodox Monastery. (χ.χ.). \emph{Papadic Formulae} [διαδικτυακή συλλογή μελωδικών τύπων από 758 συνθέσεις]. \url{https://stanthonysmonastery.org/pages/papadic-formulae}
\item Troelsgård, C. (2011). \emph{Byzantine Neumes: A New Introduction to the Middle Byzantine Musical Notation} (Monumenta Musicae Byzantinae, Subsidia 9). Copenhagen: Museum Tusculanum Press.
\end{reflist}
\subsection*{Υπολογιστική μουσικολογία και μηχανική μάθηση}
\begin{reflist}
\item Πολυκαρπίδης, Π. (2023). \emph{Υπολογιστική προσέγγιση της Βυζαντινής μουσικής: Η περίπτωση του Ειρμολογίου του Θεοφάνη Καρύκη} (Διδακτορική διατριβή). Εθνικό και Καποδιστριακό Πανεπιστήμιο Αθηνών.
\item Alexandru, M., De Luca, E., \& Fujinaga, I. (2024). Brainstorming on the encoding of Eastern music scripts, Middle-Byzantine neumes as a case-study. \emph{Series Musicologica Balcanica}, 5, 177--205. \url{https://doi.org/10.26262/smb.v0i5.10414}
\item Bauckhage, C., Schulz, D., \& von Rueden, L. (2025). Introduction and overview. In D. Schulz \& C. Bauckhage (Eds.), \emph{Informed Machine Learning} (pp. 1--14). Springer. \url{https://doi.org/10.1007/978-3-031-83097-6_1}
\item Bouris, N. G., Georgaki, A., \& Chaldaeakes, A. G. (2018). Encoding Byzantine music notation (ca. 1600--1814). \emph{Proceedings of the Sound and Music Computing Conference (SMC 2018)}.
\item Bozkurt, B., Karaosmanoğlu, M. K., Karaçalı, B., \& Ünal, E. (2014). Usul and makam driven automatic melodic segmentation for Turkish music. \emph{Journal of New Music Research}, 43(4), 375--389. \url{https://doi.org/10.1080/09298215.2014.924535}
\item Cambouropoulos, E. (2001a). The Local Boundary Detection Model (LBDM) and its application in the study of expressive timing. \emph{Proceedings of the International Computer Music Conference (ICMC 2001)}.
\item Cambouropoulos, E. (2001b). Melodic cue abstraction, similarity, and category formation: A formal model. \emph{Music Perception}, 18(3), 347--370. \url{https://doi.org/10.1525/mp.2001.18.3.347}
\item Cambouropoulos, E. (2006). Musical parallelism and melodic segmentation: A computational approach. \emph{Music Perception}, 23(3), 249--268. \url{https://doi.org/10.1525/mp.2006.23.3.249}
\item Conklin, D., \& Anagnostopoulou, C. (2006). Segmental pattern discovery in music. \emph{INFORMS Journal on Computing}, 18(3), 285--293.
\item Cornelissen, B., Zuidema, W., \& Burgoyne, J. A. (2020). Mode classification and natural units in plainchant. \emph{Proceedings of ISMIR 2020}.
\item Darwiche, A., \& Marquis, P. (2002). A knowledge compilation map. \emph{Journal of Artificial Intelligence Research}, 17, 229--264. \url{https://doi.org/10.1613/jair.989}
\item Delviniotis, D., \& Kouroupetroglou, G. (2011). DAMASKINOS: The prototype corpus of Greek Orthodox ecclesiastical chant voices. \emph{Proceedings of Crossroads: Greece as an Intercultural Pole of Musical Thought and Creativity}, Thessaloniki.
\item Finkensiep, C., Déguernel, K., Neuwirth, M., \& Rohrmeier, M. (2020). Voice-leading schema recognition using rhythm and pitch features. \emph{Proceedings of ISMIR 2020}, 520--526.
\item Gingras, B., Pearce, M. T., Goodchild, M., Dean, R. T., Wiggins, G. A., \& McAdams, S. (2016). Linking melodic expectation to expressive performance timing and perceived musical tension. \emph{Journal of Experimental Psychology: Human Perception and Performance}, 42(4), 594--609.
\item Guan, X., Dong, Z., Liu, H., \& Li, Q. (2025). Improving phrase segmentation in symbolic folk music: A hybrid model with local context and global structure awareness. \emph{Entropy}, 27(5), 460. \url{https://doi.org/10.3390/e27050460}
\item Guan, Y., Zhao, J., Qiu, Y., Zhang, Z., \& Xia, G. (2018). Melodic phrase segmentation by deep neural networks. \emph{arXiv preprint}, arXiv:1811.05688. \url{https://doi.org/10.48550/arXiv.1811.05688}
\item Holzapfel, A. (2010). \emph{Similarity Methods for Computational Ethnomusicology} (Doctoral dissertation). University of Crete. \url{https://doi.org/10.12681/eadd/30917}
\item Janssen, B., van Kranenburg, P., \& Volk, A. (2017). Finding occurrences of melodic segments in folk songs employing symbolic similarity measures. \emph{Journal of New Music Research}, 46(2), 118--134.
\item Kapoor, S., \& Narayanan, A. (2023). Leakage and the reproducibility crisis in machine-learning-based science. \emph{Patterns}, 4(9), 100804. \url{https://doi.org/10.1016/j.patter.2023.100804}
\item Karsdorp, F., van Kranenburg, P., \& Manjavacas, E. (2019). Learning similarity metrics for melody retrieval. \emph{Proceedings of ISMIR 2019}, 478--485.
\item Khosla, P., Teterwak, P., Wang, C., Sarna, A., Tian, Y., Isola, P., Maschinot, A., Liu, C., \& Krishnan, D. (2020). Supervised contrastive learning. \emph{Advances in Neural Information Processing Systems 33}, 18661--18673.
\item Knight, K., \& Graehl, J. (1998). Machine transliteration. \emph{Computational Linguistics}, 24(4), 599--612.
\item Koh, P. W., Nguyen, T., Tang, Y. S., Mussmann, S., Pierson, E., Kim, B., \& Liang, P. (2020). Concept bottleneck models. \emph{Proceedings of the 37th International Conference on Machine Learning (ICML)}, PMLR 119, 5338--5348.
\item Kong, L., Dyer, C., \& Smith, N. A. (2016). Segmental recurrent neural networks. \emph{Proceedings of the International Conference on Learning Representations (ICLR 2016)}. \url{https://arxiv.org/abs/1511.06018}
\item Lanz, V., \& Hajič Jr., J. (2025). Gregorian melody, modality, and memory: Segmenting chant with Bayesian nonparametrics. \emph{Proceedings of ISMIR 2025}. \url{https://doi.org/10.5281/zenodo.17811449}
\item Marconato, E., Teso, S., Vergari, A., \& Passerini, A. (2023). Not all neuro-symbolic concepts are created equal: Analysis and mitigation of reasoning shortcuts. \emph{Advances in Neural Information Processing Systems 36} (NeurIPS 2023).
\item Mavromatis, P. (2005a). A hidden Markov model of melody production in Greek church chant. \emph{Computing in Musicology}, 14, 93--112.
\item Mavromatis, P. (2005b). \emph{The Echoi of Modern Greek Church Chant in Written and Oral Transmission: A Computational Model and Its Cognitive Implications} (Doctoral dissertation). Eastman School of Music, University of Rochester.
\item Mongeau, M., \& Sankoff, D. (1990). Comparison of musical sequences. \emph{Computers and the Humanities}, 24(3), 161--175.
\item Neanes Project. (χ.χ.). \emph{Neanes: A Byzantine chant scorewriter} [λογισμικό ανοικτού κώδικα]. \url{https://github.com/neanes/neanes}.
\item Nuttall, T., García Casado, M., Ferraro, A., Conklin, D., \& Caro Repetto, R. (2021). A computational exploration of melodic patterns in Arab-Andalusian music. \emph{Journal of Mathematics and Music}, 15(2), 168--180. \url{https://doi.org/10.1080/17459737.2021.1917010}
\item Pearce, M. T., Müllensiefen, D., \& Wiggins, G. A. (2010). The role of expectation and probabilistic learning in auditory boundary perception: A model comparison. \emph{Perception}, 39(10), 1365--1391.
\item Pearce, M. T. (2018). Statistical learning and probabilistic prediction in music cognition: Mechanisms of stylistic enculturation. \emph{Annals of the New York Academy of Sciences}, 1423(1), 378--395.
\item Polykarpidis, P., Kalofonos, D., Balageorgos, D., \& Anagnostopoulou, C. (2022). Three related corpora in Middle Byzantine music notation and a preliminary comparative analysis. \emph{Proceedings of ISMIR 2022}, 306--313.
\item Rao, P., Ross, J. C., Ganguli, K. K., Pandit, V., Ishwar, V., Bellur, A., \& Murthy, H. A. (2014). Classification of melodic motifs in raga music with time-series matching. \emph{Journal of New Music Research}, 43(1), 115--131. \url{https://doi.org/10.1080/09298215.2013.873470}
\item Raykar, V. C., Yu, S., Zhao, L. H., Hermosillo Valadez, G., Florin, C., Bogoni, L., \& Moy, L. (2010). Learning from crowds. \emph{Journal of Machine Learning Research}, 11, 1297--1322.
\item Richard, G., Lostanlen, V., Yang, Y.-H., \& Müller, M. (2024). Model-based deep learning for music information research. \emph{IEEE Signal Processing Magazine}, 41(6), 51--59.
\item Ristad, E. S., \& Yianilos, P. N. (1998). Learning string-edit distance. \emph{IEEE Transactions on Pattern Analysis and Machine Intelligence}, 20(5), 522--532.
\item Sarawagi, S., \& Cohen, W. W. (2004). Semi-Markov conditional random fields for information extraction. \emph{Advances in Neural Information Processing Systems 17}, 1185--1192.
\item Schiødt, N. (1971). A computer-aided analysis of thirty-five Byzantine hymns. In E. Wellesz \& M. Velimirović (Eds.), \emph{Studies in Eastern Chant} (Vol. 2, pp. 129--154). London: Oxford University Press.
\item Sturm, B. L. (2014). A simple method to determine if a music information retrieval system is a ``horse''. \emph{IEEE Transactions on Multimedia}, 16(6), 1636--1644. \url{https://doi.org/10.1109/TMM.2014.2330697}
\item Thomae, M., Alexandru, M., De Luca, E., Aslanidi, M., Bazmadelis, A., Fujinaga, I., Ioannidou, A., \& Siklafidis, N. (2026, 30 June). \emph{A long-awaited desideratum: MEI encoding of Byzantine neumes (Chrysanthine notation)} [Conference presentation/abstract]. IAML Congress, Thessaloniki. \url{https://portfolio.avarts.ionio.gr/en/news/34889/}
\item Tsuboi, Y., Kashima, H., Mori, S., Oda, H., \& Matsumoto, Y. (2008). Training conditional random fields using incomplete annotations. \emph{Proceedings of COLING 2008}, 897--904.
\item van Kranenburg, P. (2020). Rule mining for local boundary detection in melodies. \emph{Proceedings of the 21st International Society for Music Information Retrieval Conference (ISMIR 2020)}.
\item van Krieken, E., Minervini, P., Ponti, E. M., \& Vergari, A. (2024). On the independence assumption in neurosymbolic learning. \emph{Proceedings of the 41st ICML}, PMLR 235, 49078--49097.
\item Volk, A., \& van Kranenburg, P. (2012). Melodic similarity among folk songs: An annotation study on similarity-based categorization in music. \emph{Musicae Scientiae}, 16(3), 317--339.
\item von Rueden, L., Mayer, S., Beckh, K., Georgiev, B., Giesselbach, S., Heese, R., Kirsch, B., Pfrommer, J., Pick, A., Ramamurthy, R., Walczak, M., Garcke, J., Bauckhage, C., \& Schuecker, J. (2023). Informed machine learning: A taxonomy and survey of integrating prior knowledge into learning systems. \emph{IEEE Transactions on Knowledge and Data Engineering}, 35(1), 614--633.
\item Wang, C., Mysore, G. J., \& Dubnov, S. (2017). Re-visiting the music segmentation problem with crowdsourcing. \emph{Proceedings of ISMIR 2017}.
\item Xu, J., Zhang, Z., Friedman, T., Liang, Y., \& Van den Broeck, G. (2018). A semantic loss function for deep learning with symbolic knowledge. \emph{Proceedings of the 35th ICML}, PMLR 80, 5502--5511.
\item Zeng, M., Tan, X., Wang, R., Ju, Z., Qin, T., \& Liu, T.-Y. (2021). MusicBERT: Symbolic music understanding with large-scale pre-training. \emph{Findings of the Association for Computational Linguistics: ACL-IJCNLP 2021}, 791--800. \url{https://doi.org/10.18653/v1/2021.findings-acl.70}
\end{reflist}
\end{document}
